% Complete reviewed Unit 17; editorial master: chapters/ch17/answer-key.md
% Markdown SHA-256: c6710364dde62f50c01f7bfeab57fd6ed813dddaa13d6860cfd55e75b5b2a81d
% One-use initial assembly followed by complete content and page review.
\newcommand{\runningtitle}{Unit 17. Answer key}
% Layout inherited from the reviewed previous course; set a running title per unit.
% This is a small hand-maintained TeX layout, not a conversion pipeline.
\RequirePackage{fix-cm}
\documentclass[11pt,twoside]{article}
\usepackage{fontspec}
\setmainfont{Cambria}[
  Renderer=HarfBuzz,
  BoldFont=cambriab.ttf,
  ItalicFont=cambriai.ttf,
  BoldItalicFont=cambriaz.ttf
]
\setmonofont{lmmono10-regular.otf}[BoldFont=lmmonolt10-bold.otf]
\usepackage[paperwidth=165mm,paperheight=241mm,inner=21mm,outer=19mm,
  top=22mm,bottom=20mm,headheight=14pt,headsep=7mm,footskip=10mm]{geometry}
\usepackage{xcolor,array,booktabs,longtable,tabularx}
\usepackage{enumitem,fancyhdr}
\usepackage[unicode,hidelinks]{hyperref}
\definecolor{ink}{HTML}{000000}
\definecolor{muted}{HTML}{4C5558}
\setlength{\parindent}{0pt}
\setlength{\parskip}{3.5pt}
\AtBeginDocument{\fontsize{11}{14.3}\selectfont}
\setlength{\tabcolsep}{4pt}
\renewcommand{\arraystretch}{1.10}
\setlist[itemize]{leftmargin=1.3em,label=\textbullet,itemsep=2pt,topsep=3pt,parsep=0pt}
\setlist[enumerate]{leftmargin=1.8em,itemsep=3pt,topsep=4pt,parsep=0pt}
\widowpenalty=10000
\clubpenalty=10000
\emergencystretch=1em
\raggedbottom
\pagestyle{fancy}
\fancyhf{}
\providecommand{\runningtitle}{Old English}
\fancyhead[LE,RO]{\small\itshape \runningtitle}
\fancyhead[LO,RE]{\small Old English}
\fancyfoot[LE,RO]{\small\thepage}
\renewcommand{\headrulewidth}{0.3pt}
\newcommand{\booktitle}[2]{
  {\fontsize{9}{11}\selectfont\color{muted}#1\par}
  \vspace{3pt}
  {\fontsize{20}{24}\selectfont\color{ink}#2\par}
  \vspace{6pt}
}
\newcommand{\source}[1]{\par{\footnotesize\color{muted}#1\par}}
\newcommand{\readingtext}[1]{
  \par\addvspace{5pt}
  \begingroup\leftskip=1.5em
    \fontsize{12}{16}\selectfont #1\par
  \endgroup\addvspace{5pt}
}
\makeatletter
\renewcommand{\section}{\@startsection{section}{1}{\z@}
  {-2.4ex \@plus -0.6ex \@minus -0.2ex}{0.8ex \@plus 0.2ex}
  {\normalfont\fontsize{12}{15}\selectfont\bfseries\color{ink}}}
\renewcommand{\subsection}{\@startsection{subsection}{2}{\z@}
  {-1.8ex \@plus -0.5ex \@minus -0.2ex}{0.5ex \@plus 0.2ex}
  {\normalfont\fontsize{11}{14}\selectfont\bfseries\color{ink}}}
\makeatother

\hyphenpenalty=10000
\exhyphenpenalty=10000
\emergencystretch=3em
\newcommand{\itempara}[2]{\par\addvspace{3pt}\noindent\hangindent=1.8em\hangafter=1\makebox[1.8em][l]{#1.}#2\par}
\newcommand{\reserve}[1]{\par\penalty-100\begingroup\dimen0=\pagegoal\advance\dimen0 by-\pagetotal\ifdim\dimen0<#1\newpage\fi\endgroup}
\begin{document}
\addfontfeatures{Ligatures=TeXOff}
\raggedright
\booktitle{UNIT 17}{Answer key}

% mdparagraph:P001:begin
The composition models are new teaching material. The reading translation renders Cook's actual Old English and was prepared independently before comparison with other authorial translations. Each model is one defensible answer: preserve the participants, essential meaning, case relations, agreement, mood and clause connections, together with the task's explicit lexical, number, omission and order controls. A grammatical alternative can miss a stated target. All seven Cook sets retain the original joint group \textbf{XII, XIII}.
% mdparagraph:P001:end

% mdparagraph:P002:begin
The supplied phrases and exact dictionary hints remain visible. Copying a given citation is not a test of independent recall; selecting the required ending, person, number or role may still be the task. The complete local numeral expressions, unusual reading frames and figurative supports are not invitations to reconstruct an untaught system unaided. English analytical answers may use equivalent precise terminology.
% mdparagraph:P002:end

% mdparagraph:P003:begin
Except where an order is specified, grammatical rearrangements that preserve the relationships are acceptable. Documented spelling and inflectional variants must be distinguished from mistakes: for example, \textbf{ðām} beside selected \textbf{ðǣm}, or a licensed feminine plural \textbf{-e} beside a selected \textbf{-a}. Such variants do not excuse wrong case, person or sense. The selected tables govern tasks that expressly ask for their forms.\textsuperscript{1}
% mdparagraph:P003:end

\reserve{7\baselineskip}
% mdparagraph:P004:begin
\section*{A. Clauses and participants}
% mdparagraph:P004:end

\reserve{8\baselineskip}
% mdparagraph:P005:begin
\subsection*{A01}
% mdparagraph:P005:end

% key: KEY-NEW-U17-A01
% mdparagraph:P006:begin
\textbf{Task:} Analyse the report. The man’s goodness is presented as simultaneous with the king’s past statement.
% mdparagraph:P006:end

% mdparagraph:P007:begin
\textbf{Input:} Se cyning sægde ðæt se mann wǣre gōd.
% mdparagraph:P007:end

% output: NEW-U17-A01-O001
\reserve{5\baselineskip}
% mdparagraph:P008:begin
\textbf{Identify the two subjects and their finite verbs.}
% mdparagraph:P008:end

% mdparagraph:P009:begin
\textbf{Se cyning} is the nominative subject of \textbf{sægde}, preterite indicative third singular. \textbf{Se mann} is the nominative subject of \textbf{wǣre}, preterite subjunctive third singular. The king speaks; the reported goodness belongs to the man.
% mdparagraph:P009:end

% output: NEW-U17-A01-O002
\reserve{5\baselineskip}
% mdparagraph:P010:begin
\textbf{Classify ðæt and explain the reported time relation.}
% mdparagraph:P010:end

% mdparagraph:P011:begin
\textbf{Ðæt} introduces the content of the statement, “that the man was good”; it is a conjunction here. The task places that goodness at the time of the king's past statement. Thus past \textbf{wǣre} does not by itself mean that the goodness preceded the statement. Calling ðæt a relative pronoun would give it an unneeded participant role.
% mdparagraph:P011:end

% output: NEW-U17-A01-O003
\reserve{5\baselineskip}
% mdparagraph:P012:begin
\textbf{Does wǣre require the conclusion that the man was not good? Explain.}
% mdparagraph:P012:end

% mdparagraph:P013:begin
No. The subjunctive marks the embedded report; it does not establish that the man was bad or that the report was false. A translation such as “The king said that the man was good” preserves both the report and the stated simultaneity. An interpretation of earlier goodness or present contrary fact misses the supplied context.
% mdparagraph:P013:end

\reserve{8\baselineskip}
% mdparagraph:P014:begin
\subsection*{A02}
% mdparagraph:P014:end

% key: KEY-NEW-U17-A02
% mdparagraph:P015:begin
\textbf{Task:} Analyse the comparison; then replace the first subject by the men while keeping the king singular.
% mdparagraph:P015:end

% mdparagraph:P016:begin
\textbf{Input:} Se mann sang bęt ðonne se cyning dyde.
% mdparagraph:P016:end

% output: NEW-U17-A02-O001
\reserve{5\baselineskip}
% mdparagraph:P017:begin
\textbf{Translate the sentence.}
% mdparagraph:P017:end

% mdparagraph:P018:begin
“The man sang better than the king did.” The man is the better singer in this comparison. “Better than the king sang” expresses the same meaning in English; reversing the two performers does not.
% mdparagraph:P018:end

% output: NEW-U17-A02-O002
\reserve{5\baselineskip}
% mdparagraph:P019:begin
\textbf{State the action repeated by dyde and explain why it is singular.}
% mdparagraph:P019:end

% mdparagraph:P020:begin
\textbf{Dyde} substitutes for the action “sang.” It is third singular because its own subject is \textbf{se cyning}, independently of the first clause's subject. It does not mean that the king performed some unspecified different action.
% mdparagraph:P020:end

% output: NEW-U17-A02-O003
\reserve{5\baselineskip}
% mdparagraph:P021:begin
\textbf{Write the complete revised Old English sentence.}
% mdparagraph:P021:end

\reserve{5\baselineskip}
% mdparagraph:P022:begin
\textbf{Ðā męnn sungon bęt ðonne se cyning dyde.}
% mdparagraph:P022:end

% mdparagraph:P023:begin
\textbf{Ðā męnn} and \textbf{sungon} replace the first singular subject and verb. The singular \textbf{se cyning dyde} stays intact. Pluralizing dyde would change the second performer; retaining sang would leave the first clause without the requested plural agreement.
% mdparagraph:P023:end

\reserve{8\baselineskip}
% mdparagraph:P024:begin
\subsection*{A03}
% mdparagraph:P024:end

% key: KEY-NEW-U17-A03
% mdparagraph:P025:begin
\textbf{Task:} Analyse the relative clause; then replace the man by the men.
% mdparagraph:P025:end

% mdparagraph:P026:begin
\textbf{Input:} Ic geseah ðone mann ðe sang.
% mdparagraph:P026:end

% output: NEW-U17-A03-O001
\reserve{5\baselineskip}
% mdparagraph:P027:begin
\textbf{Give the outer case and role of ðone mann.}
% mdparagraph:P027:end

% mdparagraph:P028:begin
\textbf{Ðone mann} is accusative singular, the direct object of \textbf{geseah}: the person whom I saw. Its role in the seeing clause determines that case.
% mdparagraph:P028:end

% output: NEW-U17-A03-O002
\reserve{5\baselineskip}
% mdparagraph:P029:begin
\textbf{Give the role represented by ðe inside the relative clause.}
% mdparagraph:P029:end

% mdparagraph:P030:begin
Inside the relative clause, \textbf{ðe} represents the singer and functions as subject of \textbf{sang}. The particle is invariant. Its antecedent's accusative case in the outer clause does not make the relative's internal role accusative.
% mdparagraph:P030:end

% output: NEW-U17-A03-O003
\reserve{5\baselineskip}
% mdparagraph:P031:begin
\textbf{Write the complete revised Old English sentence.}
% mdparagraph:P031:end

\reserve{5\baselineskip}
% mdparagraph:P032:begin
\textbf{Ic geseah ðā męnn ðe sungon.}
% mdparagraph:P032:end

% mdparagraph:P033:begin
The object becomes \textbf{ðā męnn} and the relative verb becomes \textbf{sungon}. \textbf{Ic geseah} remains first singular: the same speaker saw several men. A different defensible order is possible, but changing the speaker to a plural subject changes the task.
% mdparagraph:P033:end

\reserve{8\baselineskip}
% mdparagraph:P034:begin
\subsection*{A04}
% mdparagraph:P034:end

% key: KEY-NEW-U17-A04
% mdparagraph:P035:begin
\textbf{Task:} A concerns a possible visit tomorrow. In B the speaker treats his coming now as contrary to fact. Analyse both clauses in each sentence.
% mdparagraph:P035:end

% mdparagraph:P036:begin
\textbf{Input:} Gif hē cymð, ic fare.
% mdparagraph:P036:end

% mdparagraph:P037:begin
\textbf{Input:} Gif hē cōme, ic fōre.
% mdparagraph:P037:end

% output: NEW-U17-A04-O001
\reserve{5\baselineskip}
% mdparagraph:P038:begin
\textbf{Translate A and name the mood of each verb.}
% mdparagraph:P038:end

% mdparagraph:P039:begin
A: “If he comes, I shall go.” \textbf{Cymð} is present indicative third singular; \textbf{fare} is present indicative first singular in this stated setting. Tomorrow supplies future reference. An indicative condition is possible after gif; no separate future auxiliary is needed.
% mdparagraph:P039:end

% output: NEW-U17-A04-O002
\reserve{5\baselineskip}
% mdparagraph:P040:begin
\textbf{Translate B and name the tense and mood of each verb.}
% mdparagraph:P040:end

% mdparagraph:P041:begin
B: “If he were coming now, I would go,” or “If he came now, I would go.” \textbf{Cōme} is preterite subjunctive third singular and \textbf{fōre} preterite subjunctive first singular. Here those past forms express the stipulated contrary-to-fact present situation, not a completed past visit.
% mdparagraph:P041:end

% output: NEW-U17-A04-O003
\reserve{5\baselineskip}
% mdparagraph:P042:begin
\textbf{Explain which evidence supplies the time and factual status.}
% mdparagraph:P042:end

% mdparagraph:P043:begin
The task's possible visit tomorrow versus contrary-to-fact coming now supplies the time and factual setting. The forms distinguish the two selected constructions, but \textbf{gif} alone fixes neither a universal mood nor a time. A translation of B as “If he had come, I would have gone” changes the given time relation.
% mdparagraph:P043:end

\reserve{7\baselineskip}
% mdparagraph:P044:begin
\section*{F. Forms and grammatical roles}
% mdparagraph:P044:end

\reserve{8\baselineskip}
% mdparagraph:P045:begin
\subsection*{F01}
% mdparagraph:P045:end

% key: KEY-NEW-U17-F01
% mdparagraph:P046:begin
\textbf{Task:} Use cniht in its explicitly taught servant sense and bōc book. Translate the four expressions. In each, choose the noun's case from its role rather than from an English word ending.
% mdparagraph:P046:end

% output: NEW-U17-F01-O001
\reserve{5\baselineskip}
% mdparagraph:P047:begin
\textbf{the servant's book}
% mdparagraph:P047:end

\reserve{5\baselineskip}
% mdparagraph:P048:begin
\textbf{ðæs cnihtes bōc}
% mdparagraph:P048:end

% mdparagraph:P049:begin
\textbf{Ðæs cnihtes} is the singular genitive possessor; \textbf{bōc} is the singular head. The book does not become genitive merely because its owner is genitive. The assigned sense of cniht is servant.
% mdparagraph:P049:end

% output: NEW-U17-F01-O002
\reserve{5\baselineskip}
% mdparagraph:P050:begin
\textbf{the servants' books}
% mdparagraph:P050:end

\reserve{5\baselineskip}
% mdparagraph:P051:begin
\textbf{ðāra cnihta bēc}
% mdparagraph:P051:end

% mdparagraph:P052:begin
\textbf{Ðāra cnihta} is genitive plural and \textbf{bēc} is the plural head: several servants' books. The changed vowel in bēc belongs to the book paradigm; a regularized plural would miss the selected form.
% mdparagraph:P052:end

% output: NEW-U17-F01-O003
\reserve{5\baselineskip}
% mdparagraph:P053:begin
\textbf{I give the servant the book.}
% mdparagraph:P053:end

\reserve{5\baselineskip}
% mdparagraph:P054:begin
\textbf{Ic giefe ðǣm cnihte ðā bōc.}
% mdparagraph:P054:end

% mdparagraph:P055:begin
\textbf{Ic giefe} is first singular present; \textbf{ðǣm cnihte} is the dative recipient and \textbf{ðā bōc} the accusative thing given. Changing the order of the two objects can preserve the meaning; exchanging their case roles cannot.
% mdparagraph:P055:end

% output: NEW-U17-F01-O004
\reserve{5\baselineskip}
% mdparagraph:P056:begin
\textbf{We give the servants the books.}
% mdparagraph:P056:end

\reserve{5\baselineskip}
% mdparagraph:P057:begin
\textbf{Wē giefað ðǣm cnihtum ðā bēc.}
% mdparagraph:P057:end

% mdparagraph:P058:begin
\textbf{Wē giefað} is first plural. Both objects change number: dative \textbf{ðǣm cnihtum}, accusative \textbf{ðā bēc}. Giefað is an assertion with wē here, despite sharing its spelling with a plural command form.
% mdparagraph:P058:end

\reserve{8\baselineskip}
% mdparagraph:P059:begin
\subsection*{F02}
% mdparagraph:P059:end

% key: KEY-NEW-U17-F02
% mdparagraph:P060:begin
\textbf{Task:} Use the selected endingless accusative of cwēn queen, while retaining its genitive/\allowbreak{}dative ending. Translate the three expressions, then explain the contrast.
% mdparagraph:P060:end

% output: NEW-U17-F02-O001
\reserve{5\baselineskip}
% mdparagraph:P061:begin
\textbf{I saw the queen.}
% mdparagraph:P061:end

\reserve{5\baselineskip}
% mdparagraph:P062:begin
\textbf{Ic geseah ðā cwēn.}
% mdparagraph:P062:end

% mdparagraph:P063:begin
The queen is the accusative object of seeing: \textbf{ðā cwēn}. The endingless cwēn is the explicitly selected accusative. Documented \textbf{cwēne} can be a valid accusative alternative, but does not meet this task's endingless-form target.
% mdparagraph:P063:end

% output: NEW-U17-F02-O002
\reserve{5\baselineskip}
% mdparagraph:P064:begin
\textbf{I gave the queen the book.}
% mdparagraph:P064:end

\reserve{5\baselineskip}
% mdparagraph:P065:begin
\textbf{Ic geaf ðǣre cwēne ðā bōc.}
% mdparagraph:P065:end

% mdparagraph:P066:begin
\textbf{Ðǣre cwēne} is the dative recipient; \textbf{ðā bōc} is the accusative gift. The endingless accusative option does not license dropping the dative -e.
% mdparagraph:P066:end

% output: NEW-U17-F02-O003
\reserve{5\baselineskip}
% mdparagraph:P067:begin
\textbf{the queen's book}
% mdparagraph:P067:end

\reserve{5\baselineskip}
% mdparagraph:P068:begin
\textbf{ðǣre cwēne bōc}
% mdparagraph:P068:end

% mdparagraph:P069:begin
\textbf{Ðǣre cwēne} is genitive singular, dependent on \textbf{bōc}: the book belongs to the queen. Its spelling matches the preceding dative, while its grammatical function differs.
% mdparagraph:P069:end

% output: NEW-U17-F02-O004
\reserve{5\baselineskip}
% mdparagraph:P070:begin
\textbf{Explain what the three queen forms show about identifying case from an ending alone.}
% mdparagraph:P070:end

% mdparagraph:P071:begin
\textbf{Cwēn} is accusative in the seeing clause even without an ending; \textbf{cwēne} is dative as recipient and genitive as possessor. The governing construction and article identify the case. An unchanged ending is not proof of an unchanged grammatical role.
% mdparagraph:P071:end

\reserve{8\baselineskip}
% mdparagraph:P072:begin
\subsection*{F03}
% mdparagraph:P072:end

% key: KEY-NEW-U17-F03
% mdparagraph:P073:begin
\textbf{Task:} Use fox fox, hol hole and nest nest. Translate each sentence. Keep ordinary locative on with the dative and distinguish the two neuter plural patterns.
% mdparagraph:P073:end

% output: NEW-U17-F03-O001
\reserve{5\baselineskip}
% mdparagraph:P074:begin
\textbf{The fox is in the hole.}
% mdparagraph:P074:end

\reserve{5\baselineskip}
% mdparagraph:P075:begin
\textbf{Se fox is on ðǣm hole.}
% mdparagraph:P075:end

% mdparagraph:P076:begin
\textbf{Se fox} is nominative masculine singular. \textbf{On ðǣm hole} is a dative location, not movement into a hole. \textbf{Is} agrees with the single fox.
% mdparagraph:P076:end

% output: NEW-U17-F03-O002
\reserve{5\baselineskip}
% mdparagraph:P077:begin
\textbf{The foxes are in the holes.}
% mdparagraph:P077:end

\reserve{5\baselineskip}
% mdparagraph:P078:begin
\textbf{Ðā foxas sindon on ðǣm holum.}
% mdparagraph:P078:end

% mdparagraph:P079:begin
\textbf{Ðā foxas sindon} gives plural subject and verb; \textbf{on ðǣm holum} gives dative plural location. Changing foxas alone would leave the other requested number changes incomplete.
% mdparagraph:P079:end

% output: NEW-U17-F03-O003
\reserve{5\baselineskip}
% mdparagraph:P080:begin
\textbf{The nest is in the tree.}
% mdparagraph:P080:end

\reserve{5\baselineskip}
% mdparagraph:P081:begin
\textbf{Ðæt nest is on ðǣm trēowe.}
% mdparagraph:P081:end

% mdparagraph:P082:begin
\textbf{Ðæt nest} is the singular subject. The tree is locative dative \textbf{trēowe}, with the documented w before the ending. This is a nest located in a tree, not a nest moving toward one.
% mdparagraph:P082:end

% output: NEW-U17-F03-O004
\reserve{5\baselineskip}
% mdparagraph:P083:begin
\textbf{The nests are in the trees.}
% mdparagraph:P083:end

\reserve{5\baselineskip}
% mdparagraph:P084:begin
\textbf{Ðā nest sindon on ðǣm trēowum.}
% mdparagraph:P084:end

% mdparagraph:P085:begin
Nominative plural \textbf{nest} has no overt plural ending; \textbf{ðā} and \textbf{sindon} make its plural function clear. \textbf{Trēowum} is dative plural. An unchanged nest does not justify singular is.
% mdparagraph:P085:end

% output: NEW-U17-F03-O005
\reserve{5\baselineskip}
% mdparagraph:P086:begin
\textbf{The holes are large.}
% mdparagraph:P086:end

\reserve{5\baselineskip}
% mdparagraph:P087:begin
\textbf{Ðā holu sindon miclu.}
% mdparagraph:P087:end

% mdparagraph:P088:begin
The short neuter \textbf{hol} has nominative plural \textbf{holu}, whereas heavy \textbf{nest} has a zero ending. Predicate \textbf{miclu} is strong neuter plural. An article in the subject does not turn the separate predicate adjective weak.
% mdparagraph:P088:end

\reserve{8\baselineskip}
% mdparagraph:P089:begin
\subsection*{F04}
% mdparagraph:P089:end

% key: KEY-NEW-U17-F04
% mdparagraph:P090:begin
\textbf{Task:} Use hęrian praise. Supply only the finite form required in each stated setting, then explain the contrast. The dictionary form is visible reference help.
% mdparagraph:P090:end

% output: NEW-U17-F04-O001
\reserve{5\baselineskip}
% mdparagraph:P091:begin
\textbf{A narrator states: they praise. Give present indicative third plural.}
% mdparagraph:P091:end

\reserve{5\baselineskip}
% mdparagraph:P092:begin
\textbf{hęriað}
% mdparagraph:P092:end

% mdparagraph:P093:begin
Present indicative third plural: several people praise. This is the narrator's assertion, not a wish.
% mdparagraph:P093:end

% output: NEW-U17-F04-O002
\reserve{5\baselineskip}
% page-review repair: K17-P01: F04 imperative prompt, model and explanation
\par\begin{minipage}{\linewidth}\raggedright
% mdparagraph:P094:begin
\textbf{A speaker commands several listeners: praise! Give second-plural imperative.}
% mdparagraph:P094:end

% mdparagraph:P095:begin
\textbf{hęriað}
% mdparagraph:P095:end

% mdparagraph:P096:begin
Second-plural imperative: several listeners are commanded to praise. Its spelling coincides with the indicative answer, but its addressees and function differ.
% mdparagraph:P096:end
\end{minipage}\par


% output: NEW-U17-F04-O003
\reserve{5\baselineskip}
% mdparagraph:P097:begin
\textbf{A speaker wishes: may they praise. Give present subjunctive third plural.}
% mdparagraph:P097:end

\reserve{5\baselineskip}
% mdparagraph:P098:begin
\textbf{hęrien}
% mdparagraph:P098:end

% mdparagraph:P099:begin
Present subjunctive third plural: may they praise. The wish is expressed by the \textbf{-en} form; adding a verb meaning “be permitted” would change the construction.
% mdparagraph:P099:end

% output: NEW-U17-F04-O004
\reserve{5\baselineskip}
% mdparagraph:P100:begin
\textbf{Which answers are spelled alike, and what distinguishes their functions?}
% mdparagraph:P100:end

% mdparagraph:P101:begin
The first two answers, \textbf{hęriað}, are spelled alike. The assertion has third-plural subject reference; the command addresses the second plural. \textbf{Hęrien} has the separate subjunctive ending. The supplied infinitive identifies the verb; choosing the finite ending and analysis is the work required.
% mdparagraph:P101:end

\reserve{8\baselineskip}
% mdparagraph:P102:begin
\subsection*{F05}
% mdparagraph:P102:end

% key: KEY-NEW-U17-F05
% mdparagraph:P103:begin
\textbf{Task:} Use the explicitly taught present participles of ondrǣdan fear and hyngran be hungry. Translate the two nominal phrases, then distinguish each phrase's outer case from any object governed inside it.
% mdparagraph:P103:end

% output: NEW-U17-F05-O001
\reserve{5\baselineskip}
% mdparagraph:P104:begin
\textbf{to those fearing him: use a dative plural participle without an expressed people noun; him is the one feared}
% mdparagraph:P104:end

\reserve{5\baselineskip}
% mdparagraph:P105:begin
\textbf{hine ondrǣdendum}
% mdparagraph:P105:end

% mdparagraph:P106:begin
The nominalized participle \textbf{ondrǣdendum} is dative plural, naming the beneficiaries. \textbf{Hine} is accusative singular, the person they fear. No separate people noun or finite fear verb is required.
% mdparagraph:P106:end

% output: NEW-U17-F05-O002
\reserve{5\baselineskip}
% mdparagraph:P107:begin
\textbf{the hungry: give the weak accusative plural phrase used as a direct object}
% mdparagraph:P107:end

\reserve{5\baselineskip}
% mdparagraph:P108:begin
\textbf{ðā hyngrendan}
% mdparagraph:P108:end

% mdparagraph:P109:begin
\textbf{Ðā hyngrendan} is the article plus weak accusative plural participle: the hungry people as direct object. It does not express an impersonal finite “they hunger” construction.
% mdparagraph:P109:end

% output: NEW-U17-F05-O003
\reserve{5\baselineskip}
% mdparagraph:P110:begin
\textbf{State the two grammatical relations in the fear phrase.}
% mdparagraph:P110:end

% mdparagraph:P111:begin
The whole fear phrase is dative because of its beneficiary role. Within that phrase, \textbf{ondrǣdan} governs accusative \textbf{hine}. The dependent object does not copy the phrase's outer dative case.
% mdparagraph:P111:end

% output: NEW-U17-F05-O004
\reserve{5\baselineskip}
% mdparagraph:P112:begin
\textbf{Explain why a present participle can name people here without being a finite verb.}
% mdparagraph:P112:end

% mdparagraph:P113:begin
A participle can be used substantivally for people characterized by an action or condition. It then takes the case, number and strong/\allowbreak{}weak endings required by that noun phrase. “Present participle” names its form; it does not supply finite person agreement or an independent present-tense assertion.
% mdparagraph:P113:end

\reserve{7\baselineskip}
% mdparagraph:P114:begin
\section*{C. Composition}
% mdparagraph:P114:end

\reserve{8\baselineskip}
% mdparagraph:P115:begin
\subsection*{C01}
% mdparagraph:P115:end

% key: KEY-NEW-U17-C01
% mdparagraph:P116:begin
\textbf{Task:} Translate: The king saw the men who sang. Use an invariant relative particle.
% mdparagraph:P116:end

% output: NEW-U17-C01-O001
\reserve{5\baselineskip}
% mdparagraph:P117:begin
\textbf{Write one complete Old English sentence.}
% mdparagraph:P117:end

\reserve{5\baselineskip}
% mdparagraph:P118:begin
\textbf{Se cyning geseah ðā męnn ðe sungon.}
% mdparagraph:P118:end

% mdparagraph:P119:begin
The king is the singular subject of \textbf{geseah}; \textbf{ðā męnn} is its accusative plural object. The invariant \textbf{ðe} represents those men as the subject of plural \textbf{sungon}. A demonstrative relative can be grammatical elsewhere, but this task requests the invariant particle.
% mdparagraph:P119:end

\reserve{8\baselineskip}
% mdparagraph:P120:begin
\subsection*{C02}
% mdparagraph:P120:end

% key: KEY-NEW-U17-C02
% mdparagraph:P121:begin
\textbf{Task:} Translate: The king is as good as the men are. Express equal degree with swā … swā.
% mdparagraph:P121:end

% output: NEW-U17-C02-O001
\reserve{5\baselineskip}
% mdparagraph:P122:begin
\textbf{Write one complete Old English sentence.}
% mdparagraph:P122:end

\reserve{5\baselineskip}
% mdparagraph:P123:begin
\textbf{Se cyning is swā gōd swā ðā męnn sindon.}
% mdparagraph:P123:end

% mdparagraph:P124:begin
The two \textbf{swā} markers express equal degree: the king is as good as the men. \textbf{Gōd} is the singular predicate of the king, while \textbf{sindon} agrees with the men in the comparison clause. A comparative meaning “better than” would reverse the requested equality.
% mdparagraph:P124:end

\reserve{8\baselineskip}
% page-review repair: K17-P03 complete short C03 composition
\par\begin{minipage}{\linewidth}\raggedright
% mdparagraph:P125:begin
\subsection*{C03}
% mdparagraph:P125:end

% key: KEY-NEW-U17-C03
% mdparagraph:P126:begin
\textbf{Task:} A visit is possible tomorrow. Translate: If he comes, give (plural) him the book.
% mdparagraph:P126:end

% output: NEW-U17-C03-O001
% mdparagraph:P127:begin
\textbf{Write one complete Old English sentence.}
% mdparagraph:P127:end

% mdparagraph:P128:begin
\textbf{Gif hē cymð, giefað him ðā bōc.}
% mdparagraph:P128:end

% mdparagraph:P129:begin
\textbf{Cymð} is an indicative in the open future condition. \textbf{Giefað} addresses several listeners; \textbf{him} is the dative recipient and \textbf{ðā bōc} the accusative gift. A singular command would be valid for one listener but would miss the explicit plural target.
% mdparagraph:P129:end
\end{minipage}\par


\reserve{8\baselineskip}
% mdparagraph:P130:begin
\subsection*{C04}
% mdparagraph:P130:end

% key: KEY-NEW-U17-C04
% mdparagraph:P131:begin
\textbf{Task:} Translate: The man says that the king is good. Present the content as a report, using the present subjunctive.
% mdparagraph:P131:end

% output: NEW-U17-C04-O001
\reserve{5\baselineskip}
% mdparagraph:P132:begin
\textbf{Write one complete Old English sentence.}
% mdparagraph:P132:end

\reserve{5\baselineskip}
% mdparagraph:P133:begin
\textbf{Se mann sægð ðæt se cyning sīe gōd.}
% mdparagraph:P133:end

% mdparagraph:P134:begin
\textbf{Se mann} speaks; \textbf{se cyning} is the subject of the reported goodness. \textbf{Sīe} is present subjunctive third singular, as explicitly requested. An indicative can assert certainty in a reported clause, but would not meet this stated mood target. The report need not be false.
% mdparagraph:P134:end

\clearpage
% mdparagraph:P135:begin
\section*{Cook 37}
% mdparagraph:P135:end

% mdparagraph:P136:begin
\textbf{Cook's instruction:} [In this exercise, omit the article and the second personal pronoun throughout.]
% mdparagraph:P136:end

% mdparagraph:P137:begin
This instruction applies to all fourteen items. Retain necessary third-person pronouns such as \textbf{hīe} in the praise wishes, but omit an article and an independent second-person subject throughout. The noun catalogues in 37.8–12 continue the address; they need not be converted into fresh finite sentences. Accompanying descriptive adjectives use the selected weak address forms, while \textbf{eall} retains its separately taught strong pattern.
% mdparagraph:P137:end

\par\begin{minipage}{\linewidth}\raggedright
% mdparagraph:P138:begin
\subsection*{Cook 37.1}
% mdparagraph:P138:end

% key: KEY-EX-XII-XIII-37-01
% mdparagraph:P139:begin
Hęriað Dryhten of heofonum; hęriað hine on hēanessum.
% mdparagraph:P139:end

% output: EX-XII-XIII-37-01-O001
% mdparagraph:P140:begin
\textbf{Cook's prompt:} Praise ye the Lord from (of) the heavens; praise him in the heights.
% mdparagraph:P140:end

% mdparagraph:P141:begin
\textbf{Target and explanation:} Both \textbf{hęriað} forms are plural commands. \textbf{Dryhten} and \textbf{hine} are the persons praised, in the accusative. \textbf{Of heofonum} gives the origin and \textbf{on hēanessum} the location, both with dative plurals. Adding a second-person subject or an article would miss the whole exercise's omission rule; treating praise as an assertion would change its force.
% mdparagraph:P141:end

\end{minipage}\par\addvspace{4pt}

\par\begin{minipage}{\linewidth}\raggedright
% mdparagraph:P142:begin
\subsection*{Cook 37.2}
% mdparagraph:P142:end

% key: KEY-EX-XII-XIII-37-02
% mdparagraph:P143:begin
Hęriað hine, ealle his ęnglas; hęriað hine, eall his mægnu.
% mdparagraph:P143:end

% output: EX-XII-XIII-37-02-O001
% mdparagraph:P144:begin
\textbf{Cook's prompt:} Praise ye him, all his angels; praise him, all his powers (mægen).
% mdparagraph:P144:end

% mdparagraph:P145:begin
\textbf{Target and explanation:} The angels and powers are the groups addressed, not extra objects of praise. \textbf{Ealle ęnglas} is masculine plural; \textbf{eall mægnu} is neuter plural, with the explicitly taught sense “powers.” \textbf{His} belongs to both groups. Neither the English collective sense nor the familiar strength gloss licenses a singular reading here.
% mdparagraph:P145:end

\end{minipage}\par\addvspace{4pt}

\par\begin{minipage}{\linewidth}\raggedright
% mdparagraph:P146:begin
\subsection*{Cook 37.3}
% mdparagraph:P146:end

% key: KEY-EX-XII-XIII-37-03
% mdparagraph:P147:begin
Hęriað hine, sunne and mōna; hęriað hine, ealle steorran and lēoht.
% mdparagraph:P147:end

% output: EX-XII-XIII-37-03-O001
% mdparagraph:P148:begin
\textbf{Cook's prompt:} Praise ye him, sun and moon; praise him, all ye stars and light (lēoht).
% mdparagraph:P148:end

% mdparagraph:P149:begin
\textbf{Target and explanation:} Sun and moon are singular entities in the address; \textbf{steorran} is plural, while \textbf{lēoht} is the supplied mass noun “light.” \textbf{Ealle} modifies the stars. Keep both commands and all four named entities. English ye does not override the instruction to omit the second-person pronoun.
% mdparagraph:P149:end

\end{minipage}\par\addvspace{4pt}

\par\begin{minipage}{\linewidth}\raggedright
% mdparagraph:P150:begin
\subsection*{Cook 37.4}
% mdparagraph:P150:end

% key: KEY-EX-XII-XIII-37-04
% mdparagraph:P151:begin
Hęriað hine, heofonas heofona, and wæteru ðe ofer heofonas sindon.
% mdparagraph:P151:end

% output: EX-XII-XIII-37-04-O001
% mdparagraph:P152:begin
\textbf{Cook's prompt:} Praise him, ye heavens of heavens, and ye waters that are above (ofer, 172) the heavens.
% mdparagraph:P152:end

% mdparagraph:P153:begin
\textbf{Target and explanation:} \textbf{Heofona} is genitive plural dependent on the addressed \textbf{heofonas}. \textbf{Wæteru} supplies the plural antecedent of \textbf{ðe}, the subject of \textbf{sindon} inside the relative clause. The selected \textbf{ofer heofonas} retains the local accusative target. A genitive near the relative does not determine its internal subject role.
% mdparagraph:P153:end

\end{minipage}\par\addvspace{4pt}

\par\begin{minipage}{\linewidth}\raggedright
% mdparagraph:P154:begin
\subsection*{Cook 37.5}
% mdparagraph:P154:end

% key: KEY-EX-XII-XIII-37-05
% mdparagraph:P155:begin
Hęrien hīe naman Dryhtnes, for-ðon-ðe hē gecwæð and hīe wǣron worhte; hē bebēad and hīe wǣron gesceapene.
% mdparagraph:P155:end

% output: EX-XII-XIII-37-05-O001
% mdparagraph:P156:begin
\textbf{Cook's prompt:} Let them praise [praise they] the name of the Lord; because he spake (gecweðan), and they were made (wyrcean); he commanded, and they were created (p. 190, l. 19).
% mdparagraph:P156:end

% mdparagraph:P157:begin
\textbf{Target and explanation:} \textbf{Hęrien hīe} is a third-plural wish, “may they praise,” rather than a command addressed to them or a permission construction. \textbf{Naman} is the accusative object and \textbf{Dryhtnes} its genitive dependent. Keep the specified \textbf{gecwæð} and both passive predicates. \textbf{Worhte/\allowbreak{}gesceapene} are selected plural predicate participles after \textbf{wǣron}; worhte here is not a second finite past verb. Uninflected predicate participles are also attested, but the model illustrates the assigned plural forms.
% mdparagraph:P157:end

\end{minipage}\par\addvspace{4pt}

\par\begin{minipage}{\linewidth}\raggedright
% mdparagraph:P158:begin
\subsection*{Cook 37.6}
% mdparagraph:P158:end

% key: KEY-EX-XII-XIII-37-06
% mdparagraph:P159:begin
Hē sętte hīe ā ā ā; hē sętte bebod and ne gewīt.
% mdparagraph:P159:end

% output: EX-XII-XIII-37-06-O001
% mdparagraph:P160:begin
\textbf{Cook's prompt:} He hath appointed them for ever and ever; he (om. hath) made (sęttan) a decree (bebod, 47), and (om. it) shall not pass (gewītan).
% mdparagraph:P160:end

% mdparagraph:P161:begin
\textbf{Target and explanation:} The two \textbf{sętte} predicates appoint the entities and establish a decree. The simple preterite may render English hath appointed; the second clause expressly omits hath. The final understood subject is \textbf{bebod}, the decree, not God or plural hīe. \textbf{Ne gewīt} has future reference here. The entire compressed ending and \textbf{ā ā ā} remain supplied; restoring it would miss the explicit omission.
% mdparagraph:P161:end

\end{minipage}\par\addvspace{4pt}

\par\begin{minipage}{\linewidth}\raggedright
% mdparagraph:P162:begin
\subsection*{Cook 37.7}
% mdparagraph:P162:end

% key: KEY-EX-XII-XIII-37-07
% mdparagraph:P163:begin
Hęriað Dryhten of eorðan, dracan and ealla nēowolnessa.
% mdparagraph:P163:end

% output: EX-XII-XIII-37-07-O001
% mdparagraph:P164:begin
\textbf{Cook's prompt:} Praise the Lord from the earth, ye dragons (draca, 53), and all deeps (nēowolnes).
% mdparagraph:P164:end

% mdparagraph:P165:begin
\textbf{Target and explanation:} \textbf{Dracan} and \textbf{ealla nēowolnessa} are plural addressees; \textbf{Dryhten} is the object praised. \textbf{Of eorðan} gives the origin. The selected feminine plural quantifier is ealla. A documented -e variant need not be ungrammatical, but differs from that selected table form. Do not turn the dragons and deeps into the things being praised.
% mdparagraph:P165:end

\end{minipage}\par\addvspace{4pt}

\par\begin{minipage}{\linewidth}\raggedright
% mdparagraph:P166:begin
\subsection*{Cook 37.8}
% mdparagraph:P166:end

% key: KEY-EX-XII-XIII-37-08
% mdparagraph:P167:begin
Fȳr and hægel, snāw and īs; gāstas ȳsta ðe his word dōð;
% mdparagraph:P167:end

% output: EX-XII-XIII-37-08-O001
% mdparagraph:P168:begin
\textbf{Cook's prompt:} Fire and hail (hægel, 43); snow (snāw, 43) and ice (īs, 47); spirits of the storms (ȳst, 51. b) which fulfil (dōn) his word;
% mdparagraph:P168:end

% mdparagraph:P169:begin
\textbf{Target and explanation:} This catalogue continues the preceding address. \textbf{Ȳsta} is genitive plural, “of storms”; \textbf{ðe} represents the spirits as the subject of plural \textbf{dōð}. \textbf{His word} is their accusative object in the taught fulfilment sense of dōn. The source fragment needs neither an invented new main verb nor an added “things” noun for the listed substances.
% mdparagraph:P169:end

\end{minipage}\par\addvspace{4pt}

\par\begin{minipage}{\linewidth}\raggedright
% mdparagraph:P170:begin
\subsection*{Cook 37.9}
% mdparagraph:P170:end

% key: KEY-EX-XII-XIII-37-09
% mdparagraph:P171:begin
Muntas and ealle hyllas; wæstmbǣran trēowu and ealle cēderbēamas;
% mdparagraph:P171:end

% output: EX-XII-XIII-37-09-O001
% mdparagraph:P172:begin
\textbf{Cook's prompt:} Mountains, and all hills (hyll, 43); fruitful (wæstmbǣre, 59) trees, and all cedars;
% mdparagraph:P172:end

% mdparagraph:P173:begin
\textbf{Target and explanation:} All four groups remain plural. \textbf{Wæstmbǣran} is the selected weak adjective in direct address, agreeing with \textbf{trēowu}; the given dictionary form wæstmbǣre does not itself supply that ending. \textbf{Ealle} follows its own strong pattern with hills and cedars. Strong \textbf{wæstmbǣru} is a real neuter plural form, but does not instantiate the selected weak-address construction. Keep short \textbf{æ} in wæstm.
% mdparagraph:P173:end

\end{minipage}\par\addvspace{4pt}

\par\begin{minipage}{\linewidth}\raggedright
% mdparagraph:P174:begin
\subsection*{Cook 37.10}
% mdparagraph:P174:end

% key: KEY-EX-XII-XIII-37-10
% mdparagraph:P175:begin
Dēor and eall nīetenu; nædran and gefiðredan fuglas;
% mdparagraph:P175:end

% output: EX-XII-XIII-37-10-O001
% mdparagraph:P176:begin
\textbf{Cook's prompt:} (Wild) beasts, and all cattle; serpents, and winged (gefiðred) fowl;
% mdparagraph:P176:end

% mdparagraph:P177:begin
\textbf{Target and explanation:} \textbf{Dēor} is endingless neuter plural here; \textbf{eall nīetenu} also has neuter plural agreement. \textbf{Nædran} names the serpents. The birds are plural \textbf{fuglas}, with supplied weak address adjective \textbf{gefiðredan}. English cattle and fowl do not independently determine OE gender or allow singular agreement. The parenthetical wild does not require adding a separate adjective to dēor.
% mdparagraph:P177:end

\end{minipage}\par\addvspace{4pt}

\par\begin{minipage}{\linewidth}\raggedright
% mdparagraph:P178:begin
\subsection*{Cook 37.11}
% mdparagraph:P178:end

% key: KEY-EX-XII-XIII-37-11
% mdparagraph:P179:begin
Cyningas eorðan and eall folc; ealdormęnn and ealle dēman eorðan;
% mdparagraph:P179:end

% output: EX-XII-XIII-37-11-O001
% mdparagraph:P180:begin
\textbf{Cook's prompt:} Kings of the earth, and all peoples; princes, and all judges of the earth;
% mdparagraph:P180:end

% mdparagraph:P181:begin
\textbf{Target and explanation:} Kings, peoples, princes and judges form four addressed plural groups. \textbf{Eorðan} depends on kings and judges where it occurs. \textbf{Eall folc} is neuter plural despite the endingless noun; \textbf{ealle dēman} is masculine plural. Changing the outer case because a noun has a genitive dependent confuses the two relations.
% mdparagraph:P181:end

\end{minipage}\par\addvspace{4pt}

\par\begin{minipage}{\linewidth}\raggedright
% mdparagraph:P182:begin
\subsection*{Cook 37.12}
% mdparagraph:P182:end

% key: KEY-EX-XII-XIII-37-12
% mdparagraph:P183:begin
Geongan and fǣmnan; ealdan mid bearnum;
% mdparagraph:P183:end

% output: EX-XII-XIII-37-12-O001
% mdparagraph:P184:begin
\textbf{Cook's prompt:} Young (om. men), and maidens (fǣmne); old (om. men), with children;
% mdparagraph:P184:end

% mdparagraph:P185:begin
\textbf{Target and explanation:} \textbf{Geongan} and \textbf{ealdan} are weak plural adjectives used for the addressed young and old people. Do not restore the expressly omitted men. \textbf{Fǣmnan} is the plural noun “maidens,” while \textbf{mid bearnum} is dative plural accompaniment. Calling the adjectives substantival explains their use without creating separate new nouns.
% mdparagraph:P185:end

\end{minipage}\par\addvspace{4pt}

\par\begin{minipage}{\linewidth}\raggedright
% mdparagraph:P186:begin
\subsection*{Cook 37.13}
% mdparagraph:P186:end

% key: KEY-EX-XII-XIII-37-13
% mdparagraph:P187:begin
Hęrien hīe naman Dryhtnes; for-ðon-ðe his ānes nama is ūpāhafen; his wuldor is ofer eorðan and heofon.
% mdparagraph:P187:end

% output: EX-XII-XIII-37-13-O001
% mdparagraph:P188:begin
\textbf{Cook's prompt:} Let them praise the name of the Lord; for the name of him alone is exalted (ūpāhębban), his glory (wuldor) is above (ofer, 172) the earth and heaven.
% mdparagraph:P188:end

% mdparagraph:P189:begin
\textbf{Target and explanation:} The opening is again a third-plural wish. \textbf{His ānes} means “of him alone”; it does not mean merely “his only name.” \textbf{Ūpāhafen} is a passive participle after \textbf{is}, distinct from the finite present in 37.14. Glory is above both earth and heaven in the selected \textbf{ofer} phrase. Preserve all three predications and the causal connection.
% mdparagraph:P189:end

\end{minipage}\par\addvspace{4pt}

\par\begin{minipage}{\linewidth}\raggedright
% mdparagraph:P190:begin
\subsection*{Cook 37.14}
% mdparagraph:P190:end

% key: KEY-EX-XII-XIII-37-14
% mdparagraph:P191:begin
Hē ūpāhęfð horn his folces, lof ealra his hālgena, bearna Israeles, folces him nēalǣcendes.
% mdparagraph:P191:end

% output: EX-XII-XIII-37-14-O001
% mdparagraph:P192:begin
\textbf{Cook's prompt:} He exalteth the horn (horn, 43) of his people, the praise of all his saints (hālga, 53), of the children (bearn) of Israel (Israel), of the people drawing nigh (om. unto) him.
% mdparagraph:P192:end

% mdparagraph:P193:begin
\textbf{Target and explanation:} Retain literal \textbf{horn}. Raising it in this praise context can be understood as exalting the people's strength or power; that interpretation explains the image without establishing a replacement dictionary gloss. \textbf{Lof} is appositional praise, followed by the chain of genitives. \textbf{Hālgena} is from the explicitly cited weak noun \textbf{hālga}, not the adjective hālig. In \textbf{folces him nēalǣcendes}, the participle agrees with genitive singular folces, while him is the dative person approached. The whole dense phrase is supplied; removing the image or making him genitive loses a distinct relation.
% mdparagraph:P193:end

\end{minipage}\par\addvspace{4pt}

\clearpage
% mdparagraph:P194:begin
\section*{Cook 38}
% mdparagraph:P194:end

\par\begin{minipage}{\linewidth}\raggedright
% mdparagraph:P195:begin
\subsection*{Cook 38.1}
% mdparagraph:P195:end

% key: KEY-EX-XII-XIII-38-01
% mdparagraph:P196:begin
Se cyning worhte micel elpenbǣnen setl and oferworhte hit mid ðǣm bętstan golde.
% mdparagraph:P196:end

% output: EX-XII-XIII-38-01-O001
% mdparagraph:P197:begin
\textbf{Cook's prompt:} The king made a great ivory (elpenbǣnen, 57. 3) throne, and overlaid (oferwyrcean) it with (mid) the best gold.
% mdparagraph:P197:end

% mdparagraph:P198:begin
\textbf{Target and explanation:} \textbf{Micel elpenbǣnen setl} is the neuter accusative throne, resumed by \textbf{hit}. \textbf{Oferworhte} meets the explicit oferwyrcean target; another grammatical verb for covering would not meet that lexical instruction. \textbf{Mid ðǣm bętstan golde} gives the material, with the selected weak superlative after the article. Keep both making and overlaying.
% mdparagraph:P198:end

\end{minipage}\par\addvspace{4pt}

\par\begin{minipage}{\linewidth}\raggedright
% mdparagraph:P199:begin
\subsection*{Cook 38.2}
% mdparagraph:P199:end

% key: KEY-EX-XII-XIII-38-02
% mdparagraph:P200:begin
Sumu cwēn wæs on ðǣm dagum be sūðan, hāten Saba.
% mdparagraph:P200:end

% output: EX-XII-XIII-38-02-O001
% mdparagraph:P201:begin
\textbf{Cook's prompt:} There (om.) was a certain queen in those days in the south, called (hātan) Sheba (Saba).
% mdparagraph:P201:end

% mdparagraph:P202:begin
\textbf{Target and explanation:} The subject is feminine singular \textbf{sumu cwēn}; the existential dummy there is omitted. \textbf{On ðǣm dagum} and \textbf{be sūðan} are the supplied temporal and southern-location phrases. \textbf{Hāten Saba} identifies her name; it does not mean that she was ordered to act. The foreign name is local reference help, not an inferred productive naming paradigm.
% mdparagraph:P202:end

\end{minipage}\par\addvspace{4pt}

\par\begin{minipage}{\linewidth}\raggedright
% mdparagraph:P203:begin
\subsection*{Cook 38.3}
% mdparagraph:P203:end

% key: KEY-EX-XII-XIII-38-03
% mdparagraph:P204:begin
Ðā hīerde hēo Salomones hlīsan and cōm of ðǣm sūðernum gemǣrum tō Salomone binnan Hierusalem, and hiere olfendas bǣron sūðerne wyrta and dēorwierðe gimmas and unrīm goldes.
% mdparagraph:P204:end

% output: EX-XII-XIII-38-03-O001
% mdparagraph:P205:begin
\textbf{Cook's prompt:} Then she heard the fame (hlīsa) of Solomon (Salomones), and came from the southern (sūðerne, 59) boundaries to Solomon within Jerusalem, and her camels bore southern spices [herbs], and precious gems, and a quantity (unrīm) of gold (154. a).
% mdparagraph:P205:end

% mdparagraph:P206:begin
\textbf{Target and explanation:} Solomon possesses the fame; the queen comes from the southern boundaries to Solomon within Jerusalem. \textbf{Olfendas} is the plural camel subject of \textbf{bǣron}. Preserve herbs, gems and the quantity of gold as three coordinated loads; \textbf{goldes} is genitive after \textbf{unrīm}. The local \textbf{sūðerne wyrta} uses an evidenced feminine plural -e variant, not an error merely because another selected table uses -a. The supplied Salomones and foreign place-name remain available.
% mdparagraph:P206:end

\end{minipage}\par\addvspace{4pt}

\par\begin{minipage}{\linewidth}\raggedright
% mdparagraph:P207:begin
\subsection*{Cook 38.4}
% mdparagraph:P207:end

% key: KEY-EX-XII-XIII-38-04
% mdparagraph:P208:begin
Sēo cwēn hæfde ðā sprǣce wið Salomon and sægde him swā hwæt swā hēo on hiere heortan ðōhte.
% mdparagraph:P208:end

% output: EX-XII-XIII-38-04-O001
% mdparagraph:P209:begin
\textbf{Cook's prompt:} The queen then (ðā) had speech with (wið) Solomon, and said (om. to) him whatsoever she thought in her heart (verb last).
% mdparagraph:P209:end

% mdparagraph:P210:begin
\textbf{Target and explanation:} \textbf{Wið Salomon} is the supplied speech-partner frame, not hostility toward Solomon. \textbf{Him} is the dative person told; the content is whatever the queen thought. Keep \textbf{ðōhte} last in that subordinate clause, as directed, and retain the thought in her heart. A supported dative conversation-partner variant can be valid; it does not change the recipient/\allowbreak{}content distinction.
% mdparagraph:P210:end

\end{minipage}\par\addvspace{4pt}

\par\begin{minipage}{\linewidth}\raggedright
% mdparagraph:P211:begin
\subsection*{Cook 38.5}
% mdparagraph:P211:end

% key: KEY-EX-XII-XIII-38-05
% mdparagraph:P212:begin
Salomon hīe ðā lǣrde and sægde hiere ðæt andgiet ealra ðāra worda ðe hēo hine āscode.
% mdparagraph:P212:end

% output: EX-XII-XIII-38-05-O001
% mdparagraph:P213:begin
\textbf{Cook's prompt:} Solomon then taught her, and told (sęcgean) her the meaning of all the words that (87. d) she asked (āscian) him.
% mdparagraph:P213:end

% mdparagraph:P214:begin
\textbf{Target and explanation:} The person taught is accusative \textbf{hīe}; the person told is dative \textbf{hiere}. \textbf{Andgiet} is the meaning communicated, with dependent \textbf{ealra ðāra worda}. Inside the relative, the queen is subject and \textbf{hine} the person she asked; \textbf{ðe} represents the things asked about. Its role is not determined simply by the genitive worda outside the clause.
% mdparagraph:P214:end

\end{minipage}\par\addvspace{4pt}

\par\begin{minipage}{\linewidth}\raggedright
% mdparagraph:P215:begin
\subsection*{Cook 38.6}
% mdparagraph:P215:end

% key: KEY-EX-XII-XIII-38-06
% mdparagraph:P216:begin
Ðā geseah sēo cwēn Salomones wīsdōm and ðæt mǣre tempel ðe hē geworhte and ðæs cyninges manigfealda ðēnunga.
% mdparagraph:P216:end

% output: EX-XII-XIII-38-06-O001
% mdparagraph:P217:begin
\textbf{Cook's prompt:} Then the queen saw the wisdom of Solomon, and the great (mǣre) temple that (87. d) he (om. had) built, and the king’s manifold services.
% mdparagraph:P217:end

% mdparagraph:P218:begin
\textbf{Target and explanation:} The queen sees three things: Solomon's wisdom, the temple, and the king's services. \textbf{Mǣre} retains the named adjective. The relative's \textbf{geworhte} is simple preterite under om. had; earlier building is inferred from the context, without adding a forbidden auxiliary. \textbf{Manigfealda ðēnunga} is the selected feminine accusative plural. \textbf{Tempel} remains the full nominative/\allowbreak{}accusative stem despite syncope in other cells.
% mdparagraph:P218:end

\end{minipage}\par\addvspace{4pt}

\par\begin{minipage}{\linewidth}\raggedright
% mdparagraph:P219:begin
\subsection*{Cook 38.7}
% mdparagraph:P219:end

% key: KEY-EX-XII-XIII-38-07
% mdparagraph:P220:begin
Hēo wæs tō ðan swīðe ofwundrod, ðæt hēo nānne gāst furðor næfde.
% mdparagraph:P220:end

% output: EX-XII-XIII-38-07-O001
% mdparagraph:P221:begin
\textbf{Cook's prompt:} She was so (tō ðan) greatly astonished (ofwundrian, 118), that she had no spirit more (furðor).
% mdparagraph:P221:end

% mdparagraph:P222:begin
\textbf{Target and explanation:} The supplied \textbf{tō ðan … ðæt} relates degree of astonishment to its result. \textbf{Ofwundrod} describes the astonished subject; \textbf{nānne gāst … næfde} is one negative predication, with accusative gāst. \textbf{Furðor} preserves the “more” reference. “She was left breathless” may explain the effect, but replacing the whole literal structure with a report of her death is unsupported.
% mdparagraph:P222:end

\end{minipage}\par\addvspace{4pt}

\par\begin{minipage}{\linewidth}\raggedright
% mdparagraph:P223:begin
\subsection*{Cook 38.8}
% mdparagraph:P223:end

% key: KEY-EX-XII-XIII-38-08
% mdparagraph:P224:begin
Hēo sægde ðā tō ðǣm cyninge: Sōð is ðæt word ðe ic on mīnum ēðle be ðē and ðīnum wīsdōme hīerde; ac ic nolde ǣr gelīefan ǣr ic self hit gesāwe.
% mdparagraph:P224:end

% output: EX-XII-XIII-38-08-O001
% mdparagraph:P225:begin
\textbf{Cook's prompt:} She said then to (tō) the king, True is the word that (87. d) I heard in my country concerning thee and thy wisdom (wīsdōm); but I would not believe before I (om. my)self saw (196. b) it (verb last).
% mdparagraph:P225:end

% mdparagraph:P226:begin
\textbf{Target and explanation:} Keep both the truth of the report and the queen's former refusal to believe before seeing personally. \textbf{Self} remains: om. my does not omit it. \textbf{Gesāwe} is first-singular past subjunctive in the prospective before-clause viewed from that earlier refusal, not a claim that she never saw it. Only the final saw-clause is explicitly marked verb last; \textbf{hīerde} is also clause-final in this model, but that placement is not separately mandated. The supplied topic phrases do not make thee the object of hearing.
% mdparagraph:P226:end

\end{minipage}\par\addvspace{4pt}

\par\begin{minipage}{\linewidth}\raggedright
% mdparagraph:P227:begin
\subsection*{Cook 38.9}
% mdparagraph:P227:end

% key: KEY-EX-XII-XIII-38-09
% mdparagraph:P228:begin
Māra is ðīn wīsdōm and ðīn weorc ðonne se hlīsa wǣre ðe ic hīerde.
% mdparagraph:P228:end

% output: EX-XII-XIII-38-09-O001
% mdparagraph:P229:begin
\textbf{Cook's prompt:} Greater (micel, 66) is thy wisdom and thy work than the report (hlīsa) was (wǣre) that I heard.
% mdparagraph:P229:end

% mdparagraph:P230:begin
\textbf{Target and explanation:} \textbf{Māra} supplies the comparative, and the model treats wisdom with its work as a combined singular subject. A defensible plural matrix formulation can be grammatical if its agreement is coherent. The explicitly given \textbf{wǣre} belongs to singular \textbf{se hlīsa}, the report, not to the wisdom/\allowbreak{}work subject. Its exact spelling is visible reference help; this one supplied comparison does not require subjunctive after every ðonne.
% mdparagraph:P230:end

\end{minipage}\par\addvspace{4pt}

\par\begin{minipage}{\linewidth}\raggedright
% mdparagraph:P231:begin
\subsection*{Cook 38.10}
% mdparagraph:P231:end

% key: KEY-EX-XII-XIII-38-10
% mdparagraph:P232:begin
Ēadige sindon ðīne ðegnas and ðīne ðēowan, ðe simle ætforan ðē standað and ðīnne wīsdōm hīerað.
% mdparagraph:P232:end

% output: EX-XII-XIII-38-10-O001
% mdparagraph:P233:begin
\textbf{Cook's prompt:} Blessed (ēadig) are thy vassals and thy servants (ðēowa, 53), that always (simle) stand before (ætforan) thee, and hear thy wisdom.
% mdparagraph:P233:end

% mdparagraph:P234:begin
\textbf{Target and explanation:} Both \textbf{ðegnas} and \textbf{ðēowan} are plural subjects of \textbf{sindon}, with plural \textbf{ēadige}. They also supply the relative subject of \textbf{standað} and \textbf{hīerað}. \textbf{Ætforan ðē} locates them before the addressee; \textbf{ðīnne wīsdōm} is the masculine accusative object heard. Preserve always and both groups rather than reducing the statement to one servant.
% mdparagraph:P234:end

\end{minipage}\par\addvspace{4pt}

\clearpage
% mdparagraph:P235:begin
\section*{Cook 39}
% mdparagraph:P235:end

% mdparagraph:P236:begin
Items 39.1–4 complete the exchange concerning Solomon. Items 39.5–7 change to the building image and divine reward; they are not further events in the queen's visit.
% mdparagraph:P236:end

\par\begin{minipage}{\linewidth}\raggedright
% mdparagraph:P237:begin
\subsection*{Cook 39.1}
% mdparagraph:P237:end

% key: KEY-EX-XII-XIII-39-01
% mdparagraph:P238:begin
Blētsod sīe se ælmihtiga God, ðe ðē cēas and ðē ofer ðæt rīce Israela gesętte, ðæt ðū dōmas and rihtwīsnesse sętte.
% mdparagraph:P238:end

% output: EX-XII-XIII-39-01-O001
% mdparagraph:P239:begin
\textbf{Cook's prompt:} Blessed (blētsian) be the almighty God, who (87. d) selected thee, and placed (gesęttan) thee over (ofer, 172) the kingdom of the Israelites (Israela), that thou (om. mightest) establish (sęttan) judgments and righteousness.
% mdparagraph:P239:end

% mdparagraph:P240:begin
\textbf{Target and explanation:} \textbf{Blētsod sīe} expresses the blessing wish. God is the relative subject who chose and placed the addressee; \textbf{ðē} is accusative in those actions. The final \textbf{ðæt} introduces purpose, and \textbf{sętte} is present subjunctive second singular with \textbf{ðū}, “that you establish.” It is not a new past assertion about God. Keep \textbf{gesętte}, the accusative authority phrase after ofer, and the omission of mightest.
% mdparagraph:P240:end

\end{minipage}\par\addvspace{4pt}

\par\begin{minipage}{\linewidth}\raggedright
% mdparagraph:P241:begin
\subsection*{Cook 39.2}
% mdparagraph:P241:end

% key: KEY-EX-XII-XIII-39-02
% mdparagraph:P242:begin
Hēo forgeaf ðǣm cyninge hund and twēntig punda goldes and męnigeo dēorwierðra wyrta and dēorwierðra stāna.
% mdparagraph:P242:end

% output: EX-XII-XIII-39-02-O001
% mdparagraph:P243:begin
\textbf{Cook's prompt:} She gave (forgiefan) the king (om. a) hundred and twenty pounds (154. c) of gold (154. a), and (om. a) multitude of precious spices (wyrt) and of precious stones.
% mdparagraph:P243:end

% mdparagraph:P244:begin
\textbf{Target and explanation:} The king is dative recipient. The whole supplied amount is \textbf{hund and twēntig punda goldes}, with genitive pounds after the high number and genitive gold as material. \textbf{Męnigeo} has the dependent genitive spice and stone phrases. Retain all three gifts and omit the two articles as directed. This local 120 expression does not test a general high-numeral system.
% mdparagraph:P244:end

\end{minipage}\par\addvspace{4pt}

\par\begin{minipage}{\linewidth}\raggedright
% mdparagraph:P245:begin
\subsection*{Cook 39.3}
% mdparagraph:P245:end

% key: KEY-EX-XII-XIII-39-03
% mdparagraph:P246:begin
Salomon ēac forgeaf ðǣre cwēne eall ðæt ðe hēo æt him giernde; and hēo gewęnde ongēan tō hiere ēðle mid hiere ðegnum.
% mdparagraph:P246:end

% output: EX-XII-XIII-39-03-O001
% mdparagraph:P247:begin
\textbf{Cook's prompt:} Solomon also gave the queen whatsoever she desired (giernan) of (æt) him; and she returned (gewęndan) again (ongēan) to her country with her retainers.
% mdparagraph:P247:end

% mdparagraph:P248:begin
\textbf{Target and explanation:} The queen is dative recipient of Solomon's gift; \textbf{eall ðæt} is its accusative object. Within the relative, \textbf{giernan} has the taught desired-thing relation and \textbf{æt him} names the person asked. The invariant ðe allows the inner relation without forcing the outer object into genitive. The return keeps both her own homeland and her accompanying retainers; changing the possessor to Solomon changes the event.
% mdparagraph:P248:end

\end{minipage}\par\addvspace{4pt}

\par\begin{minipage}{\linewidth}\raggedright
% mdparagraph:P249:begin
\subsection*{Cook 39.4}
% mdparagraph:P249:end

% key: KEY-EX-XII-XIII-39-04
% mdparagraph:P250:begin
Salomon wæs ðā mǣrsod ofer eallum eorðlicum cyningum; and ealla ðēoda gewilnodon ðæt hīe hine gesāwen and his wīsdōm hīerden.
% mdparagraph:P250:end

% output: EX-XII-XIII-39-04-O001
% mdparagraph:P251:begin
\textbf{Cook's prompt:} Solomon then was magnified (mǣrsian, 118) above (ofer, 166) all earthly kings; and all nations (ðēod) desired (gewilnian) to see [that they might see, 194. b] him, and hear his wisdom.
% mdparagraph:P251:end

% mdparagraph:P252:begin
\textbf{Target and explanation:} \textbf{Ofer eallum eorðlicum cyningum} is the explicitly required dative rank comparison. The desire is expressed by the bracketed finite structure: \textbf{ðæt hīe … gesāwen … hīerden}, past subjunctive plurals after \textbf{gewilnodon}. An infinitive could express a similar meaning in an appropriate frame but misses this finite-clause target. \textbf{Ealla ðēoda} is the selected feminine plural subject.
% mdparagraph:P252:end

\end{minipage}\par\addvspace{4pt}

\par\begin{minipage}{\linewidth}\raggedright
% mdparagraph:P253:begin
\subsection*{Cook 39.5}
% mdparagraph:P253:end

% key: KEY-EX-XII-XIII-39-05
% mdparagraph:P254:begin
Swā hwā swā ofer ðissum grundwealle wyrcð gold oððe seolfor oððe dēorwierðe stānas oððe trēowu, strēaw oððe ceaf, ānes gehwilces weorc bið gesweotolod.
% mdparagraph:P254:end

% output: EX-XII-XIII-39-05-O001
% mdparagraph:P255:begin
\textbf{Cook's prompt:} Whosoever buildeth upon (ofer, 166) this foundation (grundweall, 43) gold, or silver, or precious stones, or trees, straw (strēaw, 47) or chaff (ceaf, 47), every one’s (ānes gehwilces) work shall be manifested.
% mdparagraph:P255:end

% mdparagraph:P256:begin
\textbf{Target and explanation:} Keep all six materials and the dative foundation after the specified ofer. \textbf{Swā hwā swā} is the indefinite builder; \textbf{ānes gehwilces} is the exact supplied genitive “of each one,” dependent on \textbf{weorc}. \textbf{Bið gesweotolod} expresses the future manifestation of that work. In particular, the prompt says trees: substituting a different material sense is unnecessary. Replacing the given genitive phrase with a synonym can miss a reference target.
% mdparagraph:P256:end

\end{minipage}\par\addvspace{4pt}

\par\begin{minipage}{\linewidth}\raggedright
% mdparagraph:P257:begin
\subsection*{Cook 39.6}
% mdparagraph:P257:end

% key: KEY-EX-XII-XIII-39-06
% mdparagraph:P258:begin
Gif sēo getimbrung ðurhwunað and ðǣm fȳre wiðstęnt, ðonne onfēhð se wyrhta edlēane æt Gode his weorces.
% mdparagraph:P258:end

% output: EX-XII-XIII-39-06-O001
% mdparagraph:P259:begin
\textbf{Cook's prompt:} If the building (getimbrung, 51. 3) continues (ðurhwunian, 118), and withstands the fire (164. m), then shall the workman receive (om. a) reward (edlēan, 47) from (æt) God for his work (153. e).
% mdparagraph:P259:end

% mdparagraph:P260:begin
\textbf{Target and explanation:} Feminine \textbf{sēo getimbrung} is the subject of both present indicatives; this is an open condition. \textbf{Wiðstęnt} governs dative \textbf{ðǣm fȳre}. The workman receives dative \textbf{edlēane} in the selected onfōn frame, from \textbf{æt Gode}; \textbf{his weorces} gives the causal genitive “for his work.” A documented accusative receiving frame is not universally wrong, but differs from the selected dative construction here.
% mdparagraph:P260:end

\end{minipage}\par\addvspace{4pt}

\par\begin{minipage}{\linewidth}\raggedright
% mdparagraph:P261:begin
\subsection*{Cook 39.7}
% mdparagraph:P261:end

% key: KEY-EX-XII-XIII-39-07
% mdparagraph:P262:begin
Nān ēage ne mæg on ðissum līfe gesēon, ne nān ēare hīeran, ne nānes mannes heorte āsmēan ðā ðing ðe God gearcað ðǣm ðe hine lufiað.
% mdparagraph:P262:end

% output: EX-XII-XIII-39-07-O001
% mdparagraph:P263:begin
\textbf{Cook's prompt:} No (183) eye can in this life see, nor any (183) ear (53. 2) hear, nor any (183) man’s heart conceive (āsmēan, 113), the things that God prepares (gearcian, 118) (om. for) those that love him.
% mdparagraph:P263:end

% mdparagraph:P264:begin
\textbf{Target and explanation:} Retain \textbf{mæg}: eye, ear and heart cannot see, hear or conceive the things prepared. The coordinated infinitives share that modal; \textbf{āsmēan} is a visible citation reference in a supplied frame. \textbf{Gearcað} is present, “prepares,” not a past preparation. The people who love God are dative beneficiaries without an added for. The repeated negatives express one negative scope, not cancellation.
% mdparagraph:P264:end

\end{minipage}\par\addvspace{4pt}

\clearpage
% mdparagraph:P265:begin
\section*{Cook 40}
% mdparagraph:P265:end

% mdparagraph:P266:begin
Simple preterites can render the English perfect in this passage. Where an English helper is expressly omitted, retain that target. An appropriately formed analytic perfect can be valid elsewhere; its mere resemblance to English is not required here.
% mdparagraph:P266:end

\par\begin{minipage}{\linewidth}\raggedright
% mdparagraph:P267:begin
\subsection*{Cook 40.1}
% mdparagraph:P267:end

% key: KEY-EX-XII-XIII-40-01
% mdparagraph:P268:begin
Se ęngel grētte Marian and sægde ðæt hēo mid Godes giefe āfylled wǣre, and ðæt God wǣre hiere mid; and hēo wæs blētsod betwēox wīfum.
% mdparagraph:P268:end

% output: EX-XII-XIII-40-01-O001
% mdparagraph:P269:begin
\textbf{Cook's prompt:} The angel saluted Mary, and said that she was (194. b) filled (āfyllan, 113) with God’s grace [gift], and that God was with her (preposition last); and she was blessed (blētsian) among women.
% mdparagraph:P269:end

% mdparagraph:P270:begin
\textbf{Target and explanation:} The two \textbf{wǣre} clauses report what the angel said; \textbf{hēo} is the filled person, while \textbf{God} is the subject of the second reported clause. Keep the postposed \textbf{hiere mid} as required. \textbf{Godes giefe} is the supplied grace/\allowbreak{}gift expression, not an independently inferred extension of an earlier gift gloss. The final \textbf{hēo wæs blētsod} is the separate assertion about Mary among women. The complete supplied \textbf{betwēox wīfum} gives the dative plural complement; independent recall of this government is not assessed.
% mdparagraph:P270:end

\end{minipage}\par\addvspace{4pt}

\par\begin{minipage}{\linewidth}\raggedright
% mdparagraph:P271:begin
\subsection*{Cook 40.2}
% mdparagraph:P271:end

% key: KEY-EX-XII-XIII-40-02
% mdparagraph:P272:begin
Mīn sāwol mǣrsað Dryhten, and mīn gāst blissode on Gode mīnum Hǣlende.
% mdparagraph:P272:end

% output: EX-XII-XIII-40-02-O001
% mdparagraph:P273:begin
\textbf{Cook's prompt:} My soul (om. doth) magnify (mǣrsian) the Lord, and my spirit (om. hath) rejoiced in God my Saviour.
% mdparagraph:P273:end

% mdparagraph:P274:begin
\textbf{Target and explanation:} \textbf{Mǣrsað} is present, magnifying in praise; \textbf{blissode} is preterite, rejoiced. The English helpers are omitted as directed. The whole \textbf{blissode on Gode mīnum Hǣlende} frame supplies rejoicing in God, with the dative apposition “my Saviour.” It neither describes physical enlargement of the Lord nor follows automatically from a previously known but different rejoicing preposition.
% mdparagraph:P274:end

\end{minipage}\par\addvspace{4pt}

\par\begin{minipage}{\linewidth}\raggedright
% mdparagraph:P275:begin
\subsection*{Cook 40.3}
% mdparagraph:P275:end

% key: KEY-EX-XII-XIII-40-03
% mdparagraph:P276:begin
For-ðon-ðe hē geseah his ðīnenne ēaðmōdnesse; sōðlice, heonanforð ealla cnēorissa mē ēadige sęcgað.
% mdparagraph:P276:end

% output: EX-XII-XIII-40-03-O001
% mdparagraph:P277:begin
\textbf{Cook's prompt:} For he regarded (gesēon) the humility of his handmaiden (ðīnen, 51. 5); verily (sōðlice), hence forth all generations (cnēoris) shall call (sęcgean) me blessed (ēadig).
% mdparagraph:P277:end

% mdparagraph:P278:begin
\textbf{Target and explanation:} \textbf{Geseah} has the taught regard sense, with the handmaid's humility as object. In the second clause \textbf{ealla cnēorissa} is feminine nominative plural and controls \textbf{sęcgað}. \textbf{Mē} is the person called, accusative, and \textbf{ēadige} is the agreeing feminine accusative description of Mary. The future reference does not require a future auxiliary. A saying/\allowbreak{}quotation interpretation that loses the object predicate misses “call me blessed.”
% mdparagraph:P278:end

\end{minipage}\par\addvspace{4pt}

\par\begin{minipage}{\linewidth}\raggedright
% mdparagraph:P279:begin
\subsection*{Cook 40.4}
% mdparagraph:P279:end

% key: KEY-EX-XII-XIII-40-04
% mdparagraph:P280:begin
For-ðon-ðe se ðe mihtig is dyde mē miclu ðing; and hālig is his nama.
% mdparagraph:P280:end

% output: EX-XII-XIII-40-04-O001
% mdparagraph:P281:begin
\textbf{Cook's prompt:} For he that (sē ðe) is mighty hath done (om. to) me great things; and holy is his name.
% mdparagraph:P281:end

% mdparagraph:P282:begin
\textbf{Target and explanation:} \textbf{Se ðe mihtig is} names the mighty one; \textbf{mē} is the dative beneficiary of \textbf{dyde}, without the omitted to. \textbf{Miclu ðing} is neuter accusative plural, “great things.” Do not substitute a masculine plural adjective simply because English things is plural. \textbf{Hālig} is the predicate of singular \textbf{his nama}; both predications remain.
% mdparagraph:P282:end

\end{minipage}\par\addvspace{4pt}

\par\begin{minipage}{\linewidth}\raggedright
% mdparagraph:P283:begin
\subsection*{Cook 40.5}
% mdparagraph:P283:end

% key: KEY-EX-XII-XIII-40-05
% mdparagraph:P284:begin
And his mildheortnes is hine ondrǣdendum of cnēorisse on cnēorisse.
% mdparagraph:P284:end

% output: EX-XII-XIII-40-05-O001
% mdparagraph:P285:begin
\textbf{Cook's prompt:} And his mercy is (om. on those) fearing (161) him from (of) generation to (on) generation.
% mdparagraph:P285:end

% mdparagraph:P286:begin
\textbf{Target and explanation:} \textbf{Ondrǣdendum} is a dative plural human substantive, the beneficiaries of mercy, while its internal object \textbf{hine} is accusative. Do not restore the expressly omitted on those. The complete \textbf{of cnēorisse on cnēorisse} phrase supplies the temporal range from generation to generation; its local uses are not inferred from spatial motion rules. The two case relations in the fear phrase match F05.
% mdparagraph:P286:end

\end{minipage}\par\addvspace{4pt}

\par\begin{minipage}{\linewidth}\raggedright
% mdparagraph:P287:begin
\subsection*{Cook 40.6}
% mdparagraph:P287:end

% key: KEY-EX-XII-XIII-40-06
% mdparagraph:P288:begin
Hē āwearp ðā mihtgan of setle and ūpāhōf ðā ēaðmōdan.
% mdparagraph:P288:end

% output: EX-XII-XIII-40-06-O001
% mdparagraph:P289:begin
\textbf{Cook's prompt:} He hath put down (āweorpan) the powerful from (of) (om. their) seat, and lifted up (ūpāhębban) the lowly.
% mdparagraph:P289:end

% mdparagraph:P290:begin
\textbf{Target and explanation:} \textbf{Ðā mihtgan} and \textbf{ðā ēaðmōdan} are weak accusative plural substantives, the people displaced and raised. The two strong preterites \textbf{āwearp} and \textbf{ūpāhōf} retain the reversal. \textbf{Of setle} is the source, with no possessive their under the omission control. Abstract “power” and “humility” would lose the affected people.
% mdparagraph:P290:end

\end{minipage}\par\addvspace{4pt}

\par\begin{minipage}{\linewidth}\raggedright
% mdparagraph:P291:begin
\subsection*{Cook 40.7}
% mdparagraph:P291:end

% key: KEY-EX-XII-XIII-40-07
% mdparagraph:P292:begin
Hē āfylde ðā hyngrendan mid gōdum, and ðā rīcan hē forlēt ǣmtige.
% mdparagraph:P292:end

% output: EX-XII-XIII-40-07-O001
% mdparagraph:P293:begin
\textbf{Cook's prompt:} He hath filled the hungering (om. ones) with good (om. things), and the rich (rīce) he hath let go empty.
% mdparagraph:P293:end

% mdparagraph:P294:begin
\textbf{Target and explanation:} \textbf{Ðā hyngrendan} is the hungry people as weak accusative plural. \textbf{Mid gōdum} uses the previously taught lexical noun “good things,” dative plural, without restoring the omitted things. \textbf{Ðā rīcan} is the object of \textbf{forlēt} in the let-go/\allowbreak{}dismiss sense; \textbf{ǣmtige} describes their condition. Do not turn hungry into an impersonal finite clause or make empty agree with singular hē.
% mdparagraph:P294:end

\end{minipage}\par\addvspace{4pt}

\par\begin{minipage}{\linewidth}\raggedright
% mdparagraph:P295:begin
\subsection*{Cook 40.8}
% mdparagraph:P295:end

% key: KEY-EX-XII-XIII-40-08
% mdparagraph:P296:begin
Hē onfēng his cnihte Israel and gemunde his mildheortnesse.
% mdparagraph:P296:end

% output: EX-XII-XIII-40-08-O001
% mdparagraph:P297:begin
\textbf{Cook's prompt:} He hath accepted his servant (cniht) Israel, and remembered his mercy.
% mdparagraph:P297:end

% mdparagraph:P298:begin
\textbf{Target and explanation:} \textbf{Onfēng his cnihte Israel} uses the selected dative person-acceptance frame, with cniht in the servant sense and the supplied name Israel. \textbf{Gemunde} preserves the distinct remembering action and its object mercy. The attested care/\allowbreak{}protection sense explains the reception, but “helped” need not replace the prompt's accepted. The foreign name does not license an invented general declension.
% mdparagraph:P298:end

\end{minipage}\par\addvspace{4pt}

\par\begin{minipage}{\linewidth}\raggedright
% mdparagraph:P299:begin
\subsection*{Cook 40.9}
% mdparagraph:P299:end

% key: KEY-EX-XII-XIII-40-09
% mdparagraph:P300:begin
Swā hē spræc tō ūrum fæderum, Abrahame and his sǣde, ā ā ā.
% mdparagraph:P300:end

% output: EX-XII-XIII-40-09-O001
% mdparagraph:P301:begin
\textbf{Cook's prompt:} As he spake (sprecan) to (tō) our fathers, Abraham and his seed, for ever and ever.
% mdparagraph:P301:end

% mdparagraph:P302:begin
\textbf{Target and explanation:} The comparison \textbf{swā} introduces what he spoke. \textbf{Tō} governs the fathers, Abraham and his seed; \textbf{his} refers to Abraham. The newly taught collective seed sense is offspring or descendants, not agricultural produce. Retaining “seed” preserves the image, while descendants is an acceptable explanatory translation. Keep the full perpetuity expression and do not turn the genitive possessor into the speaker.
% mdparagraph:P302:end

\end{minipage}\par\addvspace{4pt}

\par\begin{minipage}{\linewidth}\raggedright
% mdparagraph:P303:begin
\subsection*{Cook 40.10}
% mdparagraph:P303:end

% key: KEY-EX-XII-XIII-40-10
% mdparagraph:P304:begin
Maria wunode hiere mid swilce ðrīe mōnaðas and gewęnde ðā tō hiere āgnum hūse.
% mdparagraph:P304:end

% output: EX-XII-XIII-40-10-O001
% mdparagraph:P305:begin
\textbf{Cook's prompt:} Mary dwelt with her about (swilce) three months, and returned (gewęndan) then to her own house.
% mdparagraph:P305:end

% mdparagraph:P306:begin
\textbf{Target and explanation:} \textbf{Hiere mid} gives the woman with whom Mary stayed; \textbf{swilce} supplies approximation here. \textbf{Ðrīe mōnaðas} is accusative duration, not a location or a dative after mid. The destination \textbf{tō hiere āgnum hūse} is Mary's own house. The two actions are dwelling and returning; omitting own or assigning its possessor to the other woman changes the reference.
% mdparagraph:P306:end

\end{minipage}\par\addvspace{4pt}

\clearpage
% mdparagraph:P307:begin
\section*{Cook 41}
% mdparagraph:P307:end

\par\begin{minipage}{\linewidth}\raggedright
% mdparagraph:P308:begin
\subsection*{Cook 41.1}
% mdparagraph:P308:end

% key: KEY-EX-XII-XIII-41-01
% mdparagraph:P309:begin
God wæs ǣfre, and ǣfre wæs his wīsdōm of him ācęnned; and se wīsdōm is his Sunu, swā mihtig swā se Fæder.
% mdparagraph:P309:end

% output: EX-XII-XIII-41-01-O001
% mdparagraph:P310:begin
\textbf{Cook's prompt:} God was ever, and ever was his wisdom begotten of (of) him, and the wisdom is his Son, just as mighty as the Father.
% mdparagraph:P310:end

% mdparagraph:P311:begin
\textbf{Target and explanation:} Keep all three relations: God's everlasting existence, the everlasting begetting of his wisdom, and that wisdom's identity with his Son. The supplied \textbf{of him ācęnned} is this specific source relation, not a general newly taught passive-agent rule. \textbf{Swā mihtig swā} expresses equality with the Father. Nothing in the repeated ever requires inventing a chronological beginning for the begetting.
% mdparagraph:P311:end

\end{minipage}\par\addvspace{4pt}

\par\begin{minipage}{\linewidth}\raggedright
% mdparagraph:P312:begin
\subsection*{Cook 41.2}
% mdparagraph:P312:end

% key: KEY-EX-XII-XIII-41-02
% mdparagraph:P313:begin
Godes Gāst gefylð ealne ymbhwyrft middangeardes, and hē hielt eall ðing.
% mdparagraph:P313:end

% output: EX-XII-XIII-41-02-O001
% mdparagraph:P314:begin
\textbf{Cook's prompt:} The Spirit of God filleth all the circuit of the world, and he maintaineth all things.
% mdparagraph:P314:end

% mdparagraph:P315:begin
\textbf{Target and explanation:} \textbf{Gefylð} uses the fill sense, with accusative \textbf{ealne ymbhwyrft middangeardes}, the entire circuit of the world. \textbf{Hielt} is the selected present third singular of healdan in the explicitly taught sustain/\allowbreak{}maintain sense. \textbf{Eall ðing} is neuter accusative plural. Finishing a task and merely grasping objects would substitute different senses for the two predicates.
% mdparagraph:P315:end

\end{minipage}\par\addvspace{4pt}

\par\begin{minipage}{\linewidth}\raggedright
% mdparagraph:P316:begin
\subsection*{Cook 41.3}
% mdparagraph:P316:end

% key: KEY-EX-XII-XIII-41-03
% mdparagraph:P317:begin
Se ðe mīnes Fæder willan wyrcð, hē is mīn brōðor and mīn mōdor and mīn sweostor.
% mdparagraph:P317:end

% output: EX-XII-XIII-41-03-O001
% mdparagraph:P318:begin
\textbf{Cook's prompt:} He that doeth (wyrcean) my Father’s will, he is my brother, and my mother, and my sister.
% mdparagraph:P318:end

% mdparagraph:P319:begin
\textbf{Target and explanation:} \textbf{Se ðe} names whoever performs the Father's will; \textbf{hē} resumes that subject in the main clause. \textbf{Mīnes Fæder willan} is the accusative will with its genitive possessor. Keep all three predicate relations—brother, mother and sister. The figurative kinship does not require altering the grammar to claim literal biological identity.
% mdparagraph:P319:end

\end{minipage}\par\addvspace{4pt}

\par\begin{minipage}{\linewidth}\raggedright
% mdparagraph:P320:begin
\subsection*{Cook 41.4}
% mdparagraph:P320:end

% key: KEY-EX-XII-XIII-41-04
% mdparagraph:P321:begin
Wē nān ðing tō ðissum middangearde ne brōhton, ne magon wē nān ðing heonon ūs mid lǣdan.
% mdparagraph:P321:end

% output: EX-XII-XIII-41-04-O001
% mdparagraph:P322:begin
\textbf{Cook's prompt:} We brought (bringan) no thing to this world, nor can we take (lǣdan) any (183) thing hence with us.
% mdparagraph:P322:end

% mdparagraph:P323:begin
\textbf{Target and explanation:} The first clause denies bringing anything into the world; the second denies ability to carry anything away. \textbf{Brōhton} is past plural; \textbf{magon … lǣdan} is present plural modal plus infinitive. \textbf{Nān … ne} preserves negative concord, and \textbf{ūs mid} means with us. Do not replace the carrying-things sense with an unrelated leading-people interpretation or omit hence.
% mdparagraph:P323:end

\end{minipage}\par\addvspace{4pt}

\par\begin{minipage}{\linewidth}\raggedright
% mdparagraph:P324:begin
\subsection*{Cook 41.5}
% mdparagraph:P324:end

% key: KEY-EX-XII-XIII-41-05
% mdparagraph:P325:begin
Biddað, and ēow bið seald; sēcað, and gē findað; cnuciað, and ēow bið geopenod.
% mdparagraph:P325:end

% output: EX-XII-XIII-41-05-O001
% mdparagraph:P326:begin
\textbf{Cook's prompt:} Ask (biddan), and (om. it shall) be given you; seek, and ye shall find; knock (cnucian, 118), and it shall be opened (om. to) you.
% mdparagraph:P326:end

% mdparagraph:P327:begin
\textbf{Target and explanation:} Keep all three command/\allowbreak{}result pairs. \textbf{Biddað, sēcað, cnuciað} address several people; \textbf{gē findað} asserts the resulting finding with future reference. \textbf{Ēow bið seald/\allowbreak{}geopenod} are the supplied passives with dative recipients. Om. it shall does not remove the OE finite passive auxiliary: it prevents a mechanical English calque. Restoring a dummy it or a preposition before ēow misses the stated targets.
% mdparagraph:P327:end

\end{minipage}\par\addvspace{4pt}

\par\begin{minipage}{\linewidth}\raggedright
% mdparagraph:P328:begin
\subsection*{Cook 41.6}
% mdparagraph:P328:end

% key: KEY-EX-XII-XIII-41-06
% mdparagraph:P329:begin
Ūre Fæder, ðū ðe eart on heofonum, sīe ðīn nama gehālgod.
% mdparagraph:P329:end

% output: EX-XII-XIII-41-06-O001
% mdparagraph:P330:begin
\textbf{Cook's prompt:} Our Father, thou which (87. d) art in heaven(s), be thy name hallowed.
% mdparagraph:P330:end

% mdparagraph:P331:begin
\textbf{Target and explanation:} \textbf{Ūre Fæder} is direct address. The relative refers to \textbf{ðū}, so its verb is second-singular \textbf{eart}, despite English which. The name is the subject of the wish \textbf{sīe … gehālgod}. A third-person relative verb or a plain assertion that the name is holy would miss the specified person or the wish.
% mdparagraph:P331:end

\end{minipage}\par\addvspace{4pt}

\par\begin{minipage}{\linewidth}\raggedright
% mdparagraph:P332:begin
\subsection*{Cook 41.7}
% mdparagraph:P332:end

% key: KEY-EX-XII-XIII-41-07
% mdparagraph:P333:begin
Tōcume ðīn rīce; geweorðe ðīn willa on eorðan swā-swā on heofonum.
% mdparagraph:P333:end

% output: EX-XII-XIII-41-07-O001
% mdparagraph:P334:begin
\textbf{Cook's prompt:} Thy kingdom (rīce) come (tōcuman; verb first); thy will befall (geweorðan) on earth as (swā-swā) in heaven(s) (verb first).
% mdparagraph:P334:end

% mdparagraph:P335:begin
\textbf{Target and explanation:} Both clauses are wishes, with present subjunctive verbs first. \textbf{Tōcume} retains the explicitly required prefix; bare cuman may be grammatical in another formulation but does not meet that lexical target. \textbf{Geweorðe} means come to pass here. Keep will as the second subject and the supplied \textbf{swā-swā} comparison between earth and heavens.
% mdparagraph:P335:end

\end{minipage}\par\addvspace{4pt}

\par\begin{minipage}{\linewidth}\raggedright
% mdparagraph:P336:begin
\subsection*{Cook 41.8}
% mdparagraph:P336:end

% key: KEY-EX-XII-XIII-41-08
% mdparagraph:P337:begin
Sęle ūs tōdæg ūrne dæghwāmlīcan hlāf; and forgief ūs ūre gyltas, swā wē ūrum gyltendum forgiefað.
% mdparagraph:P337:end

% output: EX-XII-XIII-41-08-O001
% mdparagraph:P338:begin
\textbf{Cook's prompt:} Give (sęllan) us to-day our daily (dæghwāmlīc, 55) bread; and forgive (forgiefan) us our debts (gylt, 43), as we forgive (164) our debtors (gyltend, 43. 6).
% mdparagraph:P338:end

% mdparagraph:P339:begin
\textbf{Target and explanation:} \textbf{Sęle} and \textbf{forgief} are singular commands, addressed to God; \textbf{ūs} is dative. The daily bread is accusative \textbf{ūrne dæghwāmlīcan hlāf}. This model chooses a weak adjective after the possessive; strong \textbf{ūrne dæghwāmlīcne hlāf} is also valid and meets the prompt, which specifies neither choice. \textbf{Gyltas} are the debts, but \textbf{gyltendum} are the dative debtors forgiven in the comparison clause. \textbf{Wē … forgiefað} is first-plural assertion. Exchanging debts and debtors changes both participants and government.
% mdparagraph:P339:end

\end{minipage}\par\addvspace{4pt}

\par\begin{minipage}{\linewidth}\raggedright
% mdparagraph:P340:begin
\subsection*{Cook 41.9}
% mdparagraph:P340:end

% key: KEY-EX-XII-XIII-41-09
% mdparagraph:P341:begin
And ne lǣd ðū ūs on costnunge, ac ālīes ūs of yfele.
% mdparagraph:P341:end

% output: EX-XII-XIII-41-09-O001
% mdparagraph:P342:begin
\textbf{Cook's prompt:} And lead thou us not into (on, 172) temptation (costnung), but deliver us from evil.
% mdparagraph:P342:end

% mdparagraph:P343:begin
\textbf{Target and explanation:} \textbf{Ne lǣd ðū} is a negative singular command; \textbf{ālīes} is its positive counterpart, deliver. \textbf{Ūs} is accusative in both actions. \textbf{On costnunge} gives the directed temptation phrase, while \textbf{of yfele} gives the source of deliverance. Do not turn these into assertions about what God is doing or reverse the two prepositional relations.
% mdparagraph:P343:end

\end{minipage}\par\addvspace{4pt}

\par\begin{minipage}{\linewidth}\raggedright
% mdparagraph:P344:begin
\subsection*{Cook 41.10}
% mdparagraph:P344:end

% key: KEY-EX-XII-XIII-41-10
% mdparagraph:P345:begin
Sīe hit swā.
% mdparagraph:P345:end

% output: EX-XII-XIII-41-10-O001
% mdparagraph:P346:begin
\textbf{Cook's prompt:} Be it so.
% mdparagraph:P346:end

% mdparagraph:P347:begin
\textbf{Target and explanation:} \textbf{Sīe} is present subjunctive third singular in an independent wish or assent formula, with \textbf{hit} as subject: “May it be so.” A simple “It is so” would turn the wish into an assertion. No permission verb is needed.
% mdparagraph:P347:end

\end{minipage}\par\addvspace{4pt}

\par\begin{minipage}{\linewidth}\raggedright
% mdparagraph:P348:begin
\subsection*{Cook 41.11}
% mdparagraph:P348:end

% key: KEY-EX-XII-XIII-41-11
% mdparagraph:P349:begin
Hwilc fæder sęlð his bearne stān, gif hit hine hlāfes bitt?
% mdparagraph:P349:end

% output: EX-XII-XIII-41-11-O001
% mdparagraph:P350:begin
\textbf{Cook's prompt:} What (hwilc) father will give to his child (164. a) a stone, if it asks (156. b; 159. b) bread of him?
% mdparagraph:P350:end

% mdparagraph:P351:begin
\textbf{Target and explanation:} The father gives to dative \textbf{his bearne}. Inside the condition \textbf{hit} refers to that child and is the asker; \textbf{hine} is the father asked, accusative, and \textbf{hlāfes} the thing requested, genitive. \textbf{Bitt} agrees with the singular child. Giving and asking therefore have different case frames; English of does not require an added OE preposition.
% mdparagraph:P351:end

\end{minipage}\par\addvspace{4pt}

\par\begin{minipage}{\linewidth}\raggedright
% mdparagraph:P352:begin
\subsection*{Cook 41.12}
% mdparagraph:P352:end

% key: KEY-EX-XII-XIII-41-12
% mdparagraph:P353:begin
God āwrāt ðā ealdan ǣ mid his fingre on ðǣm stānenum tabelum.
% mdparagraph:P353:end

% output: EX-XII-XIII-41-12-O001
% mdparagraph:P354:begin
\textbf{Cook's prompt:} God wrote (āwrītan) the old law (ǣ) with his finger (finger, 43) on the stone tables (tabele, 53).
% mdparagraph:P354:end

% mdparagraph:P355:begin
\textbf{Target and explanation:} \textbf{Āwrāt} retains the prefixed write lexeme and its strong past stem. \textbf{Ðā ealdan ǣ} is the law written; \textbf{mid his fingre} is the instrument. The tablets are the dative plural surface \textbf{on ðǣm stānenum tabelum}, with stānen specifying stone material. Do not reinterpret tables as furniture or make God's finger the thing written.
% mdparagraph:P355:end

\end{minipage}\par\addvspace{4pt}

\clearpage
% mdparagraph:P356:begin
\section*{Cook 42}
% mdparagraph:P356:end

% mdparagraph:P357:begin
Keep the complete supplied frames available, especially the genitive harper phrase, the 144,000 exception, the collective company description and the blessing formula. Their unfamiliar components are not separate recall tests.
% mdparagraph:P357:end

\par\begin{minipage}{\linewidth}\raggedright
% mdparagraph:P358:begin
\subsection*{Cook 42.1}
% mdparagraph:P358:end

% key: KEY-EX-XII-XIII-42-01
% mdparagraph:P359:begin
Ðæs geongan mannes lufu befēoll on ðæt mǣden.
% mdparagraph:P359:end

% output: EX-XII-XIII-42-01-O001
% mdparagraph:P360:begin
\textbf{Cook's prompt:} The young man fell in love with the girl (p. 175, l. 20; cf. mīn lufu befielð on ðē, ‘my love falls on thee,’ ‘I fall in love with thee’).
% mdparagraph:P360:end

% mdparagraph:P361:begin
\textbf{Target and explanation:} The supplied idiom makes \textbf{lufu}, love, the grammatical subject; \textbf{ðæs geongan mannes} is its genitive possessor. \textbf{On ðæt mǣden} is the accusative person toward whom the love falls. \textbf{Befēoll} retains the exercise's be-prefix and the checked long past vowel. A translation with the man as literal falling subject misses the idiom, even if modern English says “the man fell in love.”
% mdparagraph:P361:end

\end{minipage}\par\addvspace{4pt}

\par\begin{minipage}{\linewidth}\raggedright
% mdparagraph:P362:begin
\subsection*{Cook 42.2}
% mdparagraph:P362:end

% key: KEY-EX-XII-XIII-42-02
% mdparagraph:P363:begin
Ic hīerde ðāra hearpera stefne mid hiera hearpum hearpiendra; and hīe sungon nīwne sang beforan ðǣm ðrymsetle and ðǣm fēower nīetenum; and nān mann ne meahte ðone sang leornian būton ðæt hundtēontig and fēowertig and fēower ðūsendo.
% mdparagraph:P363:end

% output: EX-XII-XIII-42-02-O001
% mdparagraph:P364:begin
\textbf{Cook's prompt:} I heard the voice of harpers harping with their harps; and they sung a new song before the throne, and before the four beasts; and no man could learn that song except the hundred and forty and four thousand.
% mdparagraph:P364:end

% mdparagraph:P365:begin
\textbf{Target and explanation:} The speaker hears; the harpers harp and sing. \textbf{Hearpera} and \textbf{hearpiendra} are genitive plural around accusative \textbf{stefne}; \textbf{hiera hearpum} is dative after mid. The new song is accusative \textbf{nīwne sang}; \textbf{ðone sang} refers back to it. Preserve the throne, four beasts and negative learning clause. The whole \textbf{būton … ðūsendo} expression supplies exactly 144,000; neither its ordering nor its unfamiliar shapes are recall tests. Do not label that whole exception dative or derive a general būton case rule from it.
% mdparagraph:P365:end

\end{minipage}\par\addvspace{4pt}

\par\begin{minipage}{\linewidth}\raggedright
% mdparagraph:P366:begin
\subsection*{Cook 42.3}
% mdparagraph:P366:end

% key: KEY-EX-XII-XIII-42-03
% mdparagraph:P367:begin
Ic geseah swā micle męnigu swā nān mann gerīman ne mæg, of eallum ðēodum and of ǣlcre mǣgðe, standende beforan Godes ðrymsetle, ealle mid hwītum gegierelum gescrȳdde and palmtwīg on hiera handum healdende.
% mdparagraph:P367:end

% output: EX-XII-XIII-42-03-O001
% mdparagraph:P368:begin
\textbf{Cook's prompt:} I saw so great (om. a) company (męnigu), as no man can number (gerīman), of all nations and of every tribe (mǣgð), standing before the throne (ðrymsetl) of God, all clothed (gescrȳdan) with white garments (gegierela), and holding palm-twigs (palmtwīg, 47) in their hands.
% mdparagraph:P368:end

% mdparagraph:P369:begin
\textbf{Target and explanation:} The seen company, \textbf{micle męnigu}, is accusative feminine singular. The subsequent people stand, are clothed, and hold palm twigs; \textbf{ealle}, plural \textbf{gescrȳdde}, and \textbf{hiera} refer to the members by sense. Keep both origin phrases and all three descriptions. \textbf{Palmtwīg} is accusative plural despite its zero ending. This complete supplied collective construction does not establish a general freedom to mix singular and plural agreement.
% mdparagraph:P369:end

\end{minipage}\par\addvspace{4pt}

\par\begin{minipage}{\linewidth}\raggedright
% mdparagraph:P370:begin
\subsection*{Cook 42.4}
% mdparagraph:P370:end

% key: KEY-EX-XII-XIII-42-04
% mdparagraph:P371:begin
And hīe sungon mid hlūdre stefne: Hǣlu sīe ūrum Gode, ðe ofer his ðrymsetle sitt.
% mdparagraph:P371:end

% output: EX-XII-XIII-42-04-O001
% mdparagraph:P372:begin
\textbf{Cook's prompt:} And they sang with (om. a) loud (hlūd) voice, Be salvation (om. to) our God which sitteth upon (ofer, 166) his throne.
% mdparagraph:P372:end

% mdparagraph:P373:begin
\textbf{Target and explanation:} The singers are plural, but the embedded praise is \textbf{Hǣlu sīe ūrum Gode}, with God as dative recipient and singular sīe. \textbf{Mid hlūdre stefne} gives the loud manner of singing. In the relative, God is the singular subject of \textbf{sitt}, and the prescribed \textbf{ofer his ðrymsetle} is dative. Keep the praise formula rather than substituting an unrelated rescue event.
% mdparagraph:P373:end

\end{minipage}\par\addvspace{4pt}

\par\begin{minipage}{\linewidth}\raggedright
% mdparagraph:P374:begin
\subsection*{Cook 42.5}
% mdparagraph:P374:end

% key: KEY-EX-XII-XIII-42-05
% mdparagraph:P375:begin
And ealle ęnglas stōdon on ymbhwyrfte his ðrymsetles and Gode āluton, ðus cweðende:
% mdparagraph:P375:end

% output: EX-XII-XIII-42-05-O001
% mdparagraph:P376:begin
\textbf{Cook's prompt:} And all (om. the) angels stood round about [in circuit of] his throne, and bowed (ālūtan, 103) to God, thus saying,
% mdparagraph:P376:end

% mdparagraph:P377:begin
\textbf{Target and explanation:} The angels stand in the circuit of the throne and bow to dative \textbf{Gode}. \textbf{Āluton} is the plural strong preterite of the intransitive bow verb, not the transitive bend verb used in 42.8. The complete supplied \textbf{ðus cweðende} accompanies their actions and introduces the following speech. It is a participle, not an additional finite verb with its own tense. Its accompanying-speech construction is not an independent recall target.
% mdparagraph:P377:end

\end{minipage}\par\addvspace{4pt}

\par\begin{minipage}{\linewidth}\raggedright
% mdparagraph:P378:begin
\subsection*{Cook 42.6}
% mdparagraph:P378:end

% key: KEY-EX-XII-XIII-42-06
% mdparagraph:P379:begin
Ūrum Gode sīe blētsung and beorhtnes, wīsdōm and ðancung, weorðscipe and stręngðu, ā ā ā.
% mdparagraph:P379:end

% output: EX-XII-XIII-42-06-O001
% mdparagraph:P380:begin
\textbf{Cook's prompt:} Be (om. to) our God blessing and brightness, wisdom and thanksgiving (ðancung), dignity and strength, for ever and ever.
% mdparagraph:P380:end

% mdparagraph:P381:begin
\textbf{Target and explanation:} Retain all six coordinated qualities: blessing, brightness, wisdom, thanksgiving, dignity and strength, together with the dative recipient and perpetuity formula. The supplied doxology has singular optative \textbf{sīe}. Its local formula does not establish a universal singular-agreement rule for every list of subjects. Omitting a near-synonym still loses a requested element.
% mdparagraph:P381:end

\end{minipage}\par\addvspace{4pt}

\par\begin{minipage}{\linewidth}\raggedright
% mdparagraph:P382:begin
\subsection*{Cook 42.7}
% mdparagraph:P382:end

% key: KEY-EX-XII-XIII-42-07
% mdparagraph:P383:begin
Foxas habbað holu, and heofones fuglas nest; sōðlice Mannes Sunu næfð hwǣr hē his hēafod āhylde.
% mdparagraph:P383:end

% output: EX-XII-XIII-42-07-O001
% mdparagraph:P384:begin
\textbf{Cook's prompt:} Foxes (fox, 43) have holes (hol, 47), and the birds of heaven nests (nest, 47); but (sōðlice) the Son of Man hath not where (hwǣr) he (om. may) lay (āhyldan, 113) his head.
% mdparagraph:P384:end

% mdparagraph:P385:begin
\textbf{Target and explanation:} \textbf{Holu} and \textbf{nest} are different neuter accusative plural patterns; the second possession clause shares the first clause's have. \textbf{Næfð hwǣr … āhylde} supplies the negative indefinite-place frame: he has no place where he might lay his head. \textbf{Āhylde} is present subjunctive third singular, without a separate may. \textbf{Hwǣr} is dependent here, not a direct question asking where he is.
% mdparagraph:P385:end

\end{minipage}\par\addvspace{4pt}

\par\begin{minipage}{\linewidth}\raggedright
% mdparagraph:P386:begin
\subsection*{Cook 42.8}
% mdparagraph:P386:end

% key: KEY-EX-XII-XIII-42-08
% mdparagraph:P387:begin
Ðā cęmpan wundon cynehelm of ðornum and sętton on his hēafod and hrēod on his swīðran, and bīegdon hiera cnēo beforan him, ðus cweðende: Hāl bēo ðū, Iudea cyning!
% mdparagraph:P387:end

% output: EX-XII-XIII-42-08-O001
% mdparagraph:P388:begin
\textbf{Cook's prompt:} The soldiers platted (windan, 104) (om. a) crown of (of) thorns (ðorn, 43), and placed on (172) his head, and a reed (hrēod, 47) in his right (om. hand), and bowed their knee before him, thus saying, Hale be thou, King of the Jews (Iudea).
% mdparagraph:P388:end

% mdparagraph:P389:begin
\textbf{Target and explanation:} The soldiers plait, place, and bend their knee; keep each action and the shared crown object where the English omits repetition. \textbf{Of ðornum} gives material; \textbf{on his hēafod} is the placement goal. The supplied \textbf{on his swīðran} names his right hand without restoring the omitted noun. \textbf{Hiera cnēo} retains the singular knee in the prompt. The complete supplied \textbf{ðus cweðende} accompanies the soldiers' actions and introduces their words; independent recall of that construction is not assessed. The final salute addresses one man as King of the Jews; it is the soldiers' speech, not the narrator's approval.
% mdparagraph:P389:end

\end{minipage}\par\addvspace{4pt}

\par\begin{minipage}{\linewidth}\raggedright
% mdparagraph:P390:begin
\subsection*{Cook 42.9}
% mdparagraph:P390:end

% key: KEY-EX-XII-XIII-42-09
% mdparagraph:P391:begin
And hīe spǣtton on hine and nōmon hrēod and bēoton his hēafod.
% mdparagraph:P391:end

% output: EX-XII-XIII-42-09-O001
% mdparagraph:P392:begin
\textbf{Cook's prompt:} And they spit (spǣtan, 113) on (on, 172) him, and took (om. a) reed and smote (bēatan) his head.
% mdparagraph:P392:end

% mdparagraph:P393:begin
\textbf{Target and explanation:} The same plural subject performs three past actions: \textbf{spǣtton}, \textbf{nōmon}, \textbf{bēoton}. \textbf{On hine} is the target of spitting; \textbf{hrēod} is the reed taken; \textbf{his hēafod} is the direct object struck. Preserve that change of object. Bēatan means beat/\allowbreak{}strike here; replacing it with a killing sense would add an outcome the prompt does not state.
% mdparagraph:P393:end

\end{minipage}\par\addvspace{4pt}

\clearpage
% mdparagraph:P394:begin
\section*{Cook 43}
% mdparagraph:P394:end

\par\begin{minipage}{\linewidth}\raggedright
% mdparagraph:P395:begin
\subsection*{Cook 43.1}
% mdparagraph:P395:end

% key: KEY-EX-XII-XIII-43-01
% mdparagraph:P396:begin
Ǣlc flǣsc is gærs, and ðæs flǣsces wuldor is swilce ðāra wyrta blōstma.
% mdparagraph:P396:end

% output: EX-XII-XIII-43-01-O001
% mdparagraph:P397:begin
\textbf{Cook's prompt:} All [each] flesh is grass, and the glory (wuldor) of the flesh is like [as if] the blossom of herbs.
% mdparagraph:P397:end

% mdparagraph:P398:begin
\textbf{Target and explanation:} \textbf{Ǣlc flǣsc is gærs} retains the each-flesh/\allowbreak{}grass image. \textbf{Ðæs flǣsces wuldor} is the glory of the flesh, compared by \textbf{swilce} with the blossom of the plants. The whole figurative comparison is supplied. “Human life is brief” can explain an implication but is not a complete translation of both clauses. See A05 for the distinction between their structures.
% mdparagraph:P398:end

\end{minipage}\par\addvspace{4pt}

\par\begin{minipage}{\linewidth}\raggedright
% mdparagraph:P399:begin
\subsection*{Cook 43.2}
% mdparagraph:P399:end

% key: KEY-EX-XII-XIII-43-02
% mdparagraph:P400:begin
Nān mann ne hīerde his stefne on strǣtum.
% mdparagraph:P400:end

% output: EX-XII-XIII-43-02-O001
% mdparagraph:P401:begin
\textbf{Cook's prompt:} No man heard (183) his voice in (om. the) streets.
% mdparagraph:P401:end

% mdparagraph:P402:begin
\textbf{Target and explanation:} \textbf{Nān mann ne hīerde} is one negation: no one heard. \textbf{His stefne} is the accusative thing heard, while \textbf{on strǣtum} is the dative place, without the expressly omitted article. The negative words do not cancel one another. A paraphrase that restores the article may convey the meaning but misses the omission target.
% mdparagraph:P402:end

\end{minipage}\par\addvspace{4pt}

\par\begin{minipage}{\linewidth}\raggedright
% mdparagraph:P403:begin
\subsection*{Cook 43.3}
% mdparagraph:P403:end

% key: KEY-EX-XII-XIII-43-03
% mdparagraph:P404:begin
Cleopa and ne geswīc ðū; āhęfe ðīne stefne swā-swā bieman.
% mdparagraph:P404:end

% output: EX-XII-XIII-43-03-O001
% mdparagraph:P405:begin
\textbf{Cook's prompt:} Cry (cleopian), and cease thou not, lift up (āhębban) thy voice like (swā-swā) (om. a) trumpet.
% mdparagraph:P405:end

% mdparagraph:P406:begin
\textbf{Target and explanation:} Keep all three singular commands: cry, do not cease, lift up the voice. \textbf{Cleopa} uses the cry-aloud sense; \textbf{ðīne stefne} is accusative. \textbf{Swā-swā bieman} is the supplied elliptical trumpet comparison. The exact comparative cue remains visible. It compares the lifted voice to the trumpet image; it does not replace voice by an actual trumpet as the command's object.
% mdparagraph:P406:end

\end{minipage}\par\addvspace{4pt}

\par\begin{minipage}{\linewidth}\raggedright
% mdparagraph:P407:begin
\subsection*{Cook 43.4}
% mdparagraph:P407:end

% key: KEY-EX-XII-XIII-43-04
% mdparagraph:P408:begin
Hlǣfdige, hwider gǣst ðū?
% mdparagraph:P408:end

% output: EX-XII-XIII-43-04-O001
% mdparagraph:P409:begin
\textbf{Cook's prompt:} Lady, whither goest thou?
% mdparagraph:P409:end

% mdparagraph:P410:begin
\textbf{Target and explanation:} \textbf{Hlǣfdige} addresses the lady; \textbf{hwider} asks direction, whither. \textbf{Gǣst ðū} is second-singular present. A static where-question would change the requested destination sense, while a plural verb would change the addressee's number.
% mdparagraph:P410:end

\end{minipage}\par\addvspace{4pt}

\par\begin{minipage}{\linewidth}\raggedright
% mdparagraph:P411:begin
\subsection*{Cook 43.5}
% mdparagraph:P411:end

% key: KEY-EX-XII-XIII-43-05
% mdparagraph:P412:begin
Ðū dēst yfele.
% mdparagraph:P412:end

% output: EX-XII-XIII-43-05-O001
% mdparagraph:P413:begin
\textbf{Cook's prompt:} Thou doest ill (yfele).
% mdparagraph:P413:end

% mdparagraph:P414:begin
\textbf{Target and explanation:} \textbf{Dēst} is second-singular present, and \textbf{yfele} is the adverb ill/\allowbreak{}badly, describing the action. The sentence says that the addressee acts badly; it is not a predicate-adjective construction requiring an adjective to agree with thou.
% mdparagraph:P414:end

\end{minipage}\par\addvspace{4pt}

\par\begin{minipage}{\linewidth}\raggedright
% mdparagraph:P415:begin
\subsection*{Cook 43.6}
% mdparagraph:P415:end

% key: KEY-EX-XII-XIII-43-06
% mdparagraph:P416:begin
Hē dyde yfele, for-ðon-ðe hē tǣlde ðā cwēn, ðæs cyninges dohtor.
% mdparagraph:P416:end

% output: EX-XII-XIII-43-06-O001
% mdparagraph:P417:begin
\textbf{Cook's prompt:} He did ill, because he censured the princess, the king’s daughter.
% mdparagraph:P417:end

% mdparagraph:P418:begin
\textbf{Target and explanation:} Both \textbf{dyde} and \textbf{tǣlde} are past singular, connected causally. \textbf{Ðā cwēn} is the selected endingless accusative with the specifically taught princess sense. \textbf{Ðæs cyninges dohtor} identifies that princess in apposition. Rendering her merely as a queen would lose the contextual sense even though the headword can bear it elsewhere.
% mdparagraph:P418:end

\end{minipage}\par\addvspace{4pt}

\par\begin{minipage}{\linewidth}\raggedright
% mdparagraph:P419:begin
\subsection*{Cook 43.7}
% mdparagraph:P419:end

% key: KEY-EX-XII-XIII-43-07
% mdparagraph:P420:begin
Hīe nǣfre gerīmcræft wel ne leornodon.
% mdparagraph:P420:end

% output: EX-XII-XIII-43-07-O001
% mdparagraph:P421:begin
\textbf{Cook's prompt:} They never learned arithmetic well.
% mdparagraph:P421:end

% mdparagraph:P422:begin
\textbf{Target and explanation:} \textbf{Leornodon} is preterite plural. \textbf{Nǣfre … ne} supplies never, and \textbf{wel} retains the quality of the learning. \textbf{Gerīmcræft} is arithmetic, not any unspecified learning. Neither the two negative expressions nor the later adverb cancels the negative scope.
% mdparagraph:P422:end

\end{minipage}\par\addvspace{4pt}

\par\begin{minipage}{\linewidth}\raggedright
% mdparagraph:P423:begin
\subsection*{Cook 43.8}
% mdparagraph:P423:end

% key: KEY-EX-XII-XIII-43-08
% mdparagraph:P424:begin
Gē witon sōna hwæt gē gīet ne witon.
% mdparagraph:P424:end

% output: EX-XII-XIII-43-08-O001
% mdparagraph:P425:begin
\textbf{Cook's prompt:} You (plur.) will soon know what you do not yet know.
% mdparagraph:P425:end

% mdparagraph:P426:begin
\textbf{Target and explanation:} Both subjects are plural \textbf{gē}, so both finite verbs are \textbf{witon}. The first present has future reference with soon; the second, with \textbf{gīet ne}, states what is not yet known. The dependent \textbf{hwæt} clause is the content of that knowledge. No future auxiliary or independently recalled contracted nytan form is required.
% mdparagraph:P426:end

\end{minipage}\par\addvspace{4pt}

\par\begin{minipage}{\linewidth}\raggedright
% mdparagraph:P427:begin
\subsection*{Cook 43.9}
% mdparagraph:P427:end

% key: KEY-EX-XII-XIII-43-09
% mdparagraph:P428:begin
Plega mid ðoðore!
% mdparagraph:P428:end

% output: EX-XII-XIII-43-09-O001
% mdparagraph:P429:begin
\textbf{Cook's prompt:} Play (with) ball!
% mdparagraph:P429:end

% mdparagraph:P430:begin
\textbf{Target and explanation:} \textbf{Plega} is the selected singular command; \textbf{mid ðoðore} is the dative ball used in the play frame. The English prompt does not fix the number of listeners, so an accurately formed plural command is also defensible with that setting. The parenthetical with is an English construction cue, not an OE dictionary form or a demand to invent a different ball noun.
% mdparagraph:P430:end

\end{minipage}\par\addvspace{4pt}

\par\begin{minipage}{\linewidth}\raggedright
% mdparagraph:P431:begin
\subsection*{Cook 43.10}
% mdparagraph:P431:end

% key: KEY-EX-XII-XIII-43-10
% mdparagraph:P432:begin
Ic nān ðing ne ete, ðēah ðe ōðre męnn eten and blīðe sīen.
% mdparagraph:P432:end

% output: EX-XII-XIII-43-10-O001
% mdparagraph:P433:begin
\textbf{Cook's prompt:} I eat nothing, although other men eat and are merry.
% mdparagraph:P433:end

% mdparagraph:P434:begin
\textbf{Target and explanation:} The speaker's \textbf{ne ete} is first-singular present indicative under nān ðing. The others have third-plural \textbf{eten} and \textbf{sīen}, present subjunctives in the admitted concession. This mood does not say that they fail to eat or be merry. The supplied \textbf{ðēah ðe} is the conjunction although; its second element is not a relative pronoun meaning who.
% mdparagraph:P434:end

\end{minipage}\par\addvspace{4pt}

\par\begin{minipage}{\linewidth}\raggedright
% mdparagraph:P435:begin
\subsection*{Cook 43.11}
% mdparagraph:P435:end

% key: KEY-EX-XII-XIII-43-11
% mdparagraph:P436:begin
Ic gewilnie ðæt ic of ēow ðæt andgiet leornie ðe gē habbað.
% mdparagraph:P436:end

% output: EX-XII-XIII-43-11-O001
% mdparagraph:P437:begin
\textbf{Cook's prompt:} I desire to learn [that I may learn] from you the knowledge that you have.
% mdparagraph:P437:end

% mdparagraph:P438:begin
\textbf{Target and explanation:} \textbf{Ic gewilnie} introduces the explicitly requested finite content clause, with first-singular subjunctive \textbf{leornie}. \textbf{Of ēow} names the plural people from whom the speaker learns. The knowledge is the antecedent of \textbf{ðe}, with \textbf{gē habbað} in the relative. An infinitive could express a desire to learn in another construction but misses the bracketed finite-clause target here.
% mdparagraph:P438:end

\end{minipage}\par\addvspace{4pt}

\par\begin{minipage}{\linewidth}\raggedright
% mdparagraph:P439:begin
\subsection*{Cook 43.12}
% mdparagraph:P439:end

% key: KEY-EX-XII-XIII-43-12
% mdparagraph:P440:begin
Cēos hwilc hiera ðū wille.
% mdparagraph:P440:end

% output: EX-XII-XIII-43-12-O001
% mdparagraph:P441:begin
\textbf{Cook's prompt:} Choose which of them you (sing.) will.
% mdparagraph:P441:end

% mdparagraph:P442:begin
\textbf{Target and explanation:} \textbf{Cēos} commands one addressee. \textbf{Hwilc} represents the choice; \textbf{hiera} is a partitive genitive, of them. \textbf{Wille} is present subjunctive second singular in the indefinite choice clause. It is not an English-style future auxiliary, and the genitive is not an adjective agreeing with the selected object.
% mdparagraph:P442:end

\end{minipage}\par\addvspace{4pt}

\par\begin{minipage}{\linewidth}\raggedright
% mdparagraph:P443:begin
\subsection*{Cook 43.13}
% mdparagraph:P443:end

% key: KEY-EX-XII-XIII-43-13
% mdparagraph:P444:begin
Binnan fēawum tīdum cōmon twēgen gelǣrede and æðelborene męnn him tōgēanes.
% mdparagraph:P444:end

% output: EX-XII-XIII-43-13-O001
% mdparagraph:P445:begin
\textbf{Cook's prompt:} Within a few hours, two educated and high-born men came toward him.
% mdparagraph:P445:end

% mdparagraph:P446:begin
\textbf{Target and explanation:} \textbf{Binnan fēawum tīdum} supplies the temporal limit. \textbf{Twēgen} agrees with masculine \textbf{męnn}; both \textbf{gelǣrede} and \textbf{æðelborene} are plural descriptions of those men. Gelǣrede is the already taught lexical adjective educated/\allowbreak{}trained. The complete supplied \textbf{him tōgēanes} expresses motion toward him, with a dative complement before the preposition; independent recall of this construction is not assessed. Keep both education and birth.
% mdparagraph:P446:end

\end{minipage}\par\addvspace{4pt}

\par\begin{minipage}{\linewidth}\raggedright
% mdparagraph:P447:begin
\subsection*{Cook 43.14}
% mdparagraph:P447:end

% key: KEY-EX-XII-XIII-43-14
% mdparagraph:P448:begin
Ic gierne ðǣre æðelan fǣmnan.
% mdparagraph:P448:end

% output: EX-XII-XIII-43-14-O001
% mdparagraph:P449:begin
\textbf{Cook's prompt:} I woo the illustrious maiden.
% mdparagraph:P449:end

% mdparagraph:P450:begin
\textbf{Target and explanation:} \textbf{Gierne} is first-singular present in the marriage-seeking sense. \textbf{Ðǣre æðelan fǣmnan} is the genitive maiden wooed, with weak adjective agreement. An accusative inferred from the modern English transitive verb would miss the taught giernan government. Familiar desire vocabulary alone would not establish this particular wooing sense.
% mdparagraph:P450:end

\end{minipage}\par\addvspace{4pt}

\par\begin{minipage}{\linewidth}\raggedright
% mdparagraph:P451:begin
\subsection*{Cook 43.15}
% mdparagraph:P451:end

% key: KEY-EX-XII-XIII-43-15
% mdparagraph:P452:begin
Wē ealle sęcgað ðæt nān mann hēr Gode swā gecwēme sīe swā hē.
% mdparagraph:P452:end

% output: EX-XII-XIII-43-15-O001
% mdparagraph:P453:begin
\textbf{Cook's prompt:} We all say that no one here is (194. b) so acceptable to God as he.
% mdparagraph:P453:end

% mdparagraph:P454:begin
\textbf{Target and explanation:} \textbf{Wē ealle} are the plural speakers. Within their report, \textbf{nān mann} is singular and \textbf{sīe} is present subjunctive third singular, as the source reference directs. The complete supplied \textbf{Gode swā gecwēme} gives dative \textbf{Gode}, the person to whom acceptable; independent recall of that adjective's government is not assessed. \textbf{Swā … swā hē} preserves the degree comparison with him. The subjunctive reports their judgment without proving that it is false; the negative stays inside the reported clause. \textbf{Nān} supplies negation here; the earlier negative-concord pattern with \textbf{ne sīe} is also valid.
% mdparagraph:P454:end

\end{minipage}\par\addvspace{4pt}

\par\begin{minipage}{\linewidth}\raggedright
% mdparagraph:P455:begin
\subsection*{Cook 43.16}
% mdparagraph:P455:end

% key: KEY-EX-XII-XIII-43-16
% mdparagraph:P456:begin
Ǣrest hē wēnde ðæt hine man slōge.
% mdparagraph:P456:end

% output: EX-XII-XIII-43-16-O001
% mdparagraph:P457:begin
\textbf{Cook's prompt:} At first he thought that he should be slain [that one should slay him].
% mdparagraph:P457:end

% mdparagraph:P458:begin
\textbf{Target and explanation:} \textbf{Hē} is the thinker, while indefinite \textbf{man} is the singular actor inside the reported clause and \textbf{hine} its accusative victim. \textbf{Slōge} is past subjunctive in the expected killing. An appropriate passive could preserve the broad meaning, but the explicit bracketed target requires this active impersonal formulation. The sentence does not assert that the killing occurred.
% mdparagraph:P458:end

\end{minipage}\par\addvspace{4pt}

\par\begin{minipage}{\linewidth}\raggedright
% mdparagraph:P459:begin
\subsection*{Cook 43.17}
% mdparagraph:P459:end

% key: KEY-EX-XII-XIII-43-17
% mdparagraph:P460:begin
Wē ealle feallað tō ðīnum fōtum.
% mdparagraph:P460:end

% output: EX-XII-XIII-43-17-O001
% mdparagraph:P461:begin
\textbf{Cook's prompt:} We all fall at (tō) your (sing.) feet.
% mdparagraph:P461:end

% mdparagraph:P462:begin
\textbf{Target and explanation:} \textbf{Wē ealle feallað} is first-plural present. \textbf{Tō ðīnum fōtum} has the dative plural feet but a singular addressee as their possessor. The contextual gesture is falling at that person's feet. A plural possessive would change the addressee, not merely agree more fully with plural fōtum.
% mdparagraph:P462:end

\end{minipage}\par\addvspace{4pt}

\par\begin{minipage}{\linewidth}\raggedright
% mdparagraph:P463:begin
\subsection*{Cook 43.18}
% mdparagraph:P463:end

% key: KEY-EX-XII-XIII-43-18
% mdparagraph:P464:begin
Ic ne gelīefe ðæt ðū mīn wīf sīe.
% mdparagraph:P464:end

% output: EX-XII-XIII-43-18-O001
% mdparagraph:P465:begin
\textbf{Cook's prompt:} I (om. do) not believe that you are (194. b) my wife.
% mdparagraph:P465:end

% mdparagraph:P466:begin
\textbf{Target and explanation:} The speaker negates believing, \textbf{ic ne gelīefe}, without translating English do. In the content clause \textbf{ðū} controls second-singular subjunctive \textbf{sīe}; its spelling also serves other persons, but its subject fixes the analysis here. \textbf{Mīn wīf} is the predicate wife. Negated belief is not itself a proof that the reported proposition is false.
% mdparagraph:P466:end

\end{minipage}\par\addvspace{4pt}

\par\begin{minipage}{\linewidth}\raggedright
% mdparagraph:P467:begin
\subsection*{Cook 43.19}
% mdparagraph:P467:end

% key: KEY-EX-XII-XIII-43-19
% mdparagraph:P468:begin
Hwǣr is mīn wer?
% mdparagraph:P468:end

% output: EX-XII-XIII-43-19-O001
% mdparagraph:P469:begin
\textbf{Cook's prompt:} Where is my husband?
% mdparagraph:P469:end

% mdparagraph:P470:begin
\textbf{Target and explanation:} \textbf{Hwǣr} asks location, with singular \textbf{is}. \textbf{Mīn wer} is the speaker's husband in the newly taught spouse sense. Whither would change a location question to one of direction; generic man alone would leave out the marital relation.
% mdparagraph:P470:end

\end{minipage}\par\addvspace{4pt}

\par\begin{minipage}{\linewidth}\raggedright
% mdparagraph:P471:begin
\subsection*{Cook 43.20}
% mdparagraph:P471:end

% key: KEY-EX-XII-XIII-43-20
% mdparagraph:P472:begin
Ic lifde endleofan gēar on Ægypta lande.
% mdparagraph:P472:end

% output: EX-XII-XIII-43-20-O001
% mdparagraph:P473:begin
\textbf{Cook's prompt:} I lived eleven years in the land of the Egyptians.
% mdparagraph:P473:end

% mdparagraph:P474:begin
\textbf{Target and explanation:} \textbf{Lifde} is preterite first singular with ic. \textbf{Endleofan gēar} is the accusative duration, eleven years. \textbf{On Ægypta lande} is dative location with the supplied genitive plural ethnonym, “in the land of the Egyptians.” The duration is not another location, and the foreign people's name is not a test of an independently reconstructed ethnic paradigm.
% mdparagraph:P474:end

\end{minipage}\par\addvspace{4pt}

\reserve{7\baselineskip}
% mdparagraph:P475:begin
\section*{A. Figurative meaning}
% mdparagraph:P475:end

\reserve{8\baselineskip}
% mdparagraph:P476:begin
\subsection*{A05}
% mdparagraph:P476:end

% key: KEY-NEW-U17-A05
% mdparagraph:P477:begin
\textbf{Task:} After translating Cook 43.1, analyse both images. Keep your complete translation beside the English prompt and the supplied vocabulary.
% mdparagraph:P477:end

% output: NEW-U17-A05-O001
\reserve{5\baselineskip}
% mdparagraph:P478:begin
\textbf{Identify the two things joined by the first copular image and explain its contextual force.}
% mdparagraph:P478:end

% mdparagraph:P479:begin
The first clause identifies \textbf{ǣlc flǣsc}, each/\allowbreak{}all flesh, with \textbf{gærs}, grass, through \textbf{is}. In context the physical image suggests human frailty and transience. Calling the clause metaphorical does not remove either noun from its translation.
% mdparagraph:P479:end

% output: NEW-U17-A05-O002
\reserve{5\baselineskip}
% mdparagraph:P480:begin
\textbf{Identify the second comparison and distinguish its structure from the first.}
% mdparagraph:P480:end

% mdparagraph:P481:begin
The second compares the flesh's \textbf{wuldor}, glory, with the plants' \textbf{blōstma}, blossom, using \textbf{swilce}. The first is a copular identification; the second marks an explicit comparison. Both flesh/\allowbreak{}grass and glory/\allowbreak{}blossom must remain. No separate hypothetical event is asserted merely by the English help “as if.”
% mdparagraph:P481:end

% output: NEW-U17-A05-O003
\reserve{5\baselineskip}
% mdparagraph:P482:begin
\textbf{Explain why Human life is brief alone is not a complete translation of 43.1.}
% mdparagraph:P482:end

% mdparagraph:P483:begin
“Human life is brief” may summarize an implication, but omits both images, the reference to glory, and their different grammatical constructions. A complete translation retains what is said and then may explain its force. Interpretive paraphrase alone does not answer the composition task.
% mdparagraph:P483:end

\reserve{8\baselineskip}
% mdparagraph:P484:begin
\subsection*{A06}
% mdparagraph:P484:end

% key: KEY-NEW-U17-A06
% mdparagraph:P485:begin
\textbf{Task:} Using Cook 37.14 and 40.9 and the printed lexical help, explain how context narrows the meanings of horn and seed. Retain the actual images in translation.
% mdparagraph:P485:end

% output: NEW-U17-A06-O001
\reserve{5\baselineskip}
% mdparagraph:P486:begin
\textbf{For 37.14, distinguish the horn image from a report about a literal animal.}
% mdparagraph:P486:end

% mdparagraph:P487:begin
In 37.14 raising the \textbf{horn} can be understood as exalting the people's strength or power in the praise context. This interprets the image while retaining literal horn in the translation; it does not independently establish a replacement dictionary gloss. The clause is not a report about a literal animal.
% mdparagraph:P487:end

% output: NEW-U17-A06-O002
\reserve{5\baselineskip}
% mdparagraph:P488:begin
\textbf{For 40.9, explain what seed denotes and identify whose seed it is.}
% mdparagraph:P488:end

% mdparagraph:P489:begin
In 40.9 \textbf{his sǣde} denotes Abraham's offspring or descendants collectively. The possessor is Abraham, not the speaker or God. “His seed” preserves the image; “his descendants” explains the newly taught collective sense. This analysis task keeps the lexical help available.
% mdparagraph:P489:end

% output: NEW-U17-A06-O003
\reserve{5\baselineskip}
% mdparagraph:P490:begin
\textbf{State what the familiar inflection alone cannot establish in these two examples.}
% mdparagraph:P490:end

% mdparagraph:P491:begin
Recognizing a noun's case or paradigm does not determine its contextual sense. The same familiar shape could occur in a literal horn or agricultural seed context. Here context guides interpretation of the horn image, while seed requires its newly taught collective meaning and the correct possessor. The task's help remains available; an ending alone establishes none of these choices.
% mdparagraph:P491:end

\clearpage
% mdparagraph:P492:begin
\section*{R. The music performance}
% mdparagraph:P492:end

\reserve{8\baselineskip}
% mdparagraph:P493:begin
\subsection*{R01}
% mdparagraph:P493:end

% key: KEY-NEW-U17-R01
% mdparagraph:P494:begin
\textbf{Task:} Translate the complete assigned reading in the three stated divisions, keeping all local help available.
% mdparagraph:P494:end

% output: NEW-U17-R01-O001
\reserve{5\baselineskip}
% mdparagraph:P495:begin
\textbf{Translate 174.7–174.13, through tǣlst.}
% mdparagraph:P495:end

% mdparagraph:P496:begin
Then she went out and ordered her harp to be fetched. And as soon as she began to play the harp, she mingled its sound with pleasant singing. Then all the men began to praise her for her musical skill; Apollonius alone remained silent.
% mdparagraph:P496:end

% mdparagraph:P497:begin
Then the king said, “Apollonius, now you are behaving badly, because all the men praise my daughter for her musical skill, and you alone, by remaining silent, criticize her.”
% mdparagraph:P497:end

% mdparagraph:P498:begin
\textbf{Points to retain:} she orders the harp fetched; she combines song and harp; Apollonius alone is silent; the king treats that silence as criticism. “Had her harp fetched” is equivalent. Saying that she fetched it herself changes the causative wording.
% mdparagraph:P498:end

% output: NEW-U17-R01-O002
\reserve{5\baselineskip}
% mdparagraph:P499:begin
\textbf{Translate 174.14–175.2, through wel gelǣred.}
% mdparagraph:P499:end

% mdparagraph:P500:begin
Apollonius said, “O good king, if you permit me, I will say that I understand that your daughter has indeed chanced upon music, but has not learned it well. Now order the harp to be given to me, and then you will know what you do not yet know.”
% mdparagraph:P500:end

% mdparagraph:P501:begin
King Arcestrates said, “Apollonius, I recognize indeed that you are well trained in all things.”
% mdparagraph:P501:end

% mdparagraph:P502:begin
\textbf{Points to retain:} \textbf{gelīefst} grants permission to speak. \textbf{Mē} in the following harp command is the recipient, not the commanded giver. The daughter has chanced upon music but has not mastered it. “Then you will know what you do not yet know” preserves the future consequence; no contrary-to-fact condition is asserted.
% mdparagraph:P502:end

% output: NEW-U17-R01-O003
\reserve{5\baselineskip}
% mdparagraph:P503:begin
\textbf{Translate 175.3–175.17, through the final forðtēah.}
% mdparagraph:P503:end

% mdparagraph:P504:begin
Then the king ordered the harp to be given to Apollonius. Apollonius then went out, dressed himself, put a royal crown on his head, took the harp in his hand, and went in. He stood in such a way that the king and all those sitting around supposed that he was not Apollonius, but Apollo, the god of the heathens. Then stillness and silence fell within the hall.
% mdparagraph:P504:end

% mdparagraph:P505:begin
Apollonius took his plectrum and began skillfully to set the harp strings in motion. He mingled the harp's sound with pleasant singing. The king himself and everyone who was present cried out loudly and praised him. After this Apollonius set down the harp and acted, presenting many fine performances there that were unfamiliar and unusual to the people. Each of the things he performed greatly pleased them all.
% mdparagraph:P505:end

% mdparagraph:P506:begin
\textbf{Points to retain:} the king has the harp handed to Apollonius; Apollonius dresses and performs; the audience supposes him to be Apollo. The final verbs describe setting down the harp, acting, and presenting performances. “Things” may replace “performances” if their context is clear. The source's singular \textbf{ungecnāwen wæs} is retained in analysis even where idiomatic English uses a plural relative.
% mdparagraph:P506:end

\reserve{8\baselineskip}
% mdparagraph:P507:begin
\subsection*{R02}
% mdparagraph:P507:end

% key: KEY-NEW-U17-R02
% mdparagraph:P508:begin
\textbf{Task:} Identify each referent and the grammatical evidence. Use the supplied source notes; say when context is needed.
% mdparagraph:P508:end

% output: NEW-U17-R02-O001
\reserve{5\baselineskip}
% mdparagraph:P509:begin
\textbf{Give the referent of hēo in 174.7.}
% mdparagraph:P509:end

% mdparagraph:P510:begin
\textbf{Hēo} is the king's daughter. Its feminine singular form is compatible with that referent, but the preceding context identifies the particular woman. The form alone does not name her.
% mdparagraph:P510:end

% output: NEW-U17-R02-O002
\reserve{5\baselineskip}
% mdparagraph:P511:begin
\textbf{Give the referent of hine in 174.16 and its role.}
% mdparagraph:P511:end

% mdparagraph:P512:begin
\textbf{Hine} refers to masculine \textbf{swēgcræft}, musical art. It is the accusative object of \textbf{geleornod} in the negative perfect: she has not learned it well. Masculine grammatical gender does not make the referent the king or Apollonius.
% mdparagraph:P512:end

% output: NEW-U17-R02-O003
\reserve{5\baselineskip}
% mdparagraph:P513:begin
\textbf{Give the referent of hē in the reported clauses 175.7–8.}
% mdparagraph:P513:end

% mdparagraph:P514:begin
Both instances of \textbf{hē} refer to Apollonius, the man whom the audience sees. The clauses report their supposed identification: not Apollonius, but Apollo. They do not change the actual referent into a god.
% mdparagraph:P514:end

% output: NEW-U17-R02-O004
\reserve{5\baselineskip}
% mdparagraph:P515:begin
\textbf{Give the referent and case of him eallum in 175.16–17.}
% mdparagraph:P515:end

% mdparagraph:P516:begin
\textbf{Him eallum} refers to all the listeners and is dative plural, the people pleased by the performances. Context supplies the group; the dative construction with \textbf{līcode} supplies their experiencer role. They are not its grammatical subject.
% mdparagraph:P516:end

\reserve{8\baselineskip}
% mdparagraph:P517:begin
\subsection*{R03}
% mdparagraph:P517:end

% key: KEY-NEW-U17-R03
% mdparagraph:P518:begin
\textbf{Task:} Identify the clause relationship in each cited expression and distinguish it from a similar construction elsewhere.
% mdparagraph:P518:end

% output: NEW-U17-R03-O001
\reserve{5\baselineskip}
% mdparagraph:P519:begin
\textbf{Classify sōna swā 174.8 and explain its time relation.}
% mdparagraph:P519:end

% mdparagraph:P520:begin
\textbf{Sōna swā} introduces “as soon as.” The start of the harp playing is immediately accompanied by the blending of harp sound and song. Here swā is part of the supplied temporal expression, not an assertion that two performers are equally skilled.
% mdparagraph:P520:end

% output: NEW-U17-R03-O002
\reserve{5\baselineskip}
% mdparagraph:P521:begin
\textbf{Classify gif 174.14–15 in its context.}
% mdparagraph:P521:end

% mdparagraph:P522:begin
\textbf{Gif ðū mē gelīefst} is an open permission condition: “if you permit me [to speak].” \textbf{Gelīefst} is present indicative second singular; the following \textbf{ic sęcge} states what the speaker will say. It is neither a request to believe a proposition nor a condition assumed contrary to fact.
% mdparagraph:P522:end

% output: NEW-U17-R03-O003
\reserve{5\baselineskip}
% mdparagraph:P523:begin
\textbf{Classify swā … ðæt 175.6–7.}
% mdparagraph:P523:end

% mdparagraph:P524:begin
In \textbf{swā stōd ðæt}, the manner or degree of his stance/\allowbreak{}appearance leads to a result: the audience takes him for Apollo. This \textbf{swā … ðæt} is not the equal-degree \textbf{swā … swā} of C02 or the temporal sōna swā at the reading's beginning.
% mdparagraph:P524:end

% output: NEW-U17-R03-O004
\reserve{5\baselineskip}
% mdparagraph:P525:begin
\textbf{Explain why the two ðæt clauses 175.7–8 do not assert that Apollonius actually is Apollo.}
% mdparagraph:P525:end

% mdparagraph:P526:begin
The two clauses are the contents of \textbf{wēndon}, “they supposed”: that he was not Apollonius and that he was Apollo. \textbf{Nǣre/\allowbreak{}wǣre} are past subjunctives in that report. The stated thought explains the mistaken identification; the narrator does not assert an actual transformation.
% mdparagraph:P526:end

\reserve{8\baselineskip}
% mdparagraph:P527:begin
\subsection*{R04}
% mdparagraph:P527:end

% key: KEY-NEW-U17-R04
% mdparagraph:P528:begin
\textbf{Task:} Analyse the supplied participial and predicate expressions without changing the source spelling.
% mdparagraph:P528:end

% output: NEW-U17-R04-O001
\reserve{5\baselineskip}
% mdparagraph:P529:begin
\textbf{Explain swīgende 174.13 using original note 12a.}
% mdparagraph:P529:end

% mdparagraph:P530:begin
Original note 12a says “For \textbf{swīgiende}.” Printed \textbf{swīgende} is a present participle: remaining silent accompanies the addressee's censure. It is not an independent finite verb. Preserve the edition's swīgende in the source and identify the supplied expansion in analysis.
% mdparagraph:P530:end

% output: NEW-U17-R04-O002
\reserve{5\baselineskip}
% mdparagraph:P531:begin
\textbf{Explain the referent of ðā ymbsittendan 175.7.}
% mdparagraph:P531:end

% mdparagraph:P532:begin
\textbf{Ðā ymbsittendan} means the people sitting around, the spectators or bystanders. It is a weak nominalized present-participle expression. Together with the king, they supply the plural subject of \textbf{wēndon}; no actual people noun need be restored to the text.
% mdparagraph:P532:end

% output: NEW-U17-R04-O003
\reserve{5\baselineskip}
% mdparagraph:P533:begin
\textbf{Account for singular ungecnāwen wæs and ungewunelic 175.16.}
% mdparagraph:P533:end

% mdparagraph:P534:begin
The source construes the display of \textbf{fela fægerra ðinga} collectively in the following singular predicate \textbf{ungecnāwen wæs and ungewunelic}. Retain that wording and its supplied local explanation. It does not license a universal rule making every plural subject singular, nor does a smoother English plural authorize silently changing the Old English.
% mdparagraph:P534:end

% output: NEW-U17-R04-O004
\reserve{5\baselineskip}
% mdparagraph:P535:begin
\textbf{Identify the subject of līcode 175.17 and distinguish it from him eallum.}
% mdparagraph:P535:end

% mdparagraph:P536:begin
The subject is \textbf{ǣlc ðāra ðinga}, each of those things, hence singular \textbf{līcode}. \textbf{Him eallum} is dative plural, the people to whom each performance was pleasing. Its plurality does not control the verb.
% mdparagraph:P536:end

\reserve{8\baselineskip}
% mdparagraph:P537:begin
\subsection*{R05}
% mdparagraph:P537:end

% key: KEY-NEW-U17-R05
% mdparagraph:P538:begin
\textbf{Task:} Separate literal wording from the contextual performance sense in each expression, using the local help and original notes.
% mdparagraph:P538:end

% output: NEW-U17-R05-O001
\reserve{5\baselineskip}
% mdparagraph:P539:begin
\textbf{Explain gefēoll on swēgcræft 174.16.}
% mdparagraph:P539:end

% mdparagraph:P540:begin
The image is of falling upon musical art; original note 15 gives the contextual translation “has chanced.” Here the daughter has chanced upon or taken up music without learning it well. Treating \textbf{on swēgcræft} as a physical place into which she fell would miss the supplied idiom.
% mdparagraph:P540:end

% output: NEW-U17-R05-O002
\reserve{5\baselineskip}
% mdparagraph:P541:begin
\textbf{Explain plegode and forðtēah 175.15.}
% mdparagraph:P541:end

% mdparagraph:P542:begin
\textbf{Plegode} means acted or performed; \textbf{forðtēah} means presented the performances. Original note 6 connects the passage with comic and tragic acting, while note 8 supports the presentation sense. “Played” can work in a clearly theatrical English rendering; playing a ball game or literally drawing objects forward does not express this scene. The Latin note is explanatory evidence, not extra wording to insert into the OE translation.
% mdparagraph:P542:end

% output: NEW-U17-R05-O003
\reserve{5\baselineskip}
% page-review repair: K17-P02: keep short Apollo prompt and answer together
\par\begin{minipage}{\linewidth}\raggedright
% mdparagraph:P543:begin
\textbf{Explain the function of Apollines 175.8.}
% mdparagraph:P543:end

% mdparagraph:P544:begin
\textbf{Apollines} identifies Apollo, as original note 2 states. It is the name used in what the spectators suppose. Preserve that reported identification and the source form; do not mistake it for a second occurrence of Apollonius's name or claim that the man really became Apollo.
% mdparagraph:P544:end
\end{minipage}\par


\reserve{8\baselineskip}
% mdparagraph:P545:begin
\subsection*{R06}
% mdparagraph:P545:end

% key: KEY-NEW-U17-R06
% mdparagraph:P546:begin
\textbf{Task:} Explain the argument of the episode from the actual Old English.
% mdparagraph:P546:end

% output: NEW-U17-R06-O001
\reserve{5\baselineskip}
% mdparagraph:P547:begin
\textbf{State the difference between the king’s reading of silence and Apollonius’s explanation.}
% mdparagraph:P547:end

% mdparagraph:P548:begin
The king treats Apollonius's silence as blame and therefore as bad conduct toward his daughter. Apollonius explains a technical judgment: she has encountered music but has not learned it well. He seeks permission to speak and offers a demonstration, rather than denying that his silence conveys criticism.
% mdparagraph:P548:end

% output: NEW-U17-R06-O002
\reserve{5\baselineskip}
% mdparagraph:P549:begin
\textbf{Identify two kinds of evidence by which Apollonius answers the criticism of his conduct.}
% mdparagraph:P549:end

% mdparagraph:P550:begin
He demonstrates skilled harp playing joined with pleasant song, then displays further striking performances after putting down the harp. The listeners respond by shouting praise and taking pleasure in what he presents. Equivalent answers can describe the two kinds of performance in their own words; a bare claim that he is talented does not identify the evidence in the passage.
% mdparagraph:P550:end

% output: NEW-U17-R06-O003
\reserve{5\baselineskip}
% mdparagraph:P551:begin
\textbf{Explain how the final sentence combines novelty with approval.}
% mdparagraph:P551:end

% mdparagraph:P552:begin
The performances are unfamiliar and unusual to the people, yet each greatly pleases all of them. Novelty and approval coexist. “Unknown” does not mean bad or displeasing here, and the final dative experiencers identify the approving audience.
% mdparagraph:P552:end

% mdparagraph:P553:begin
\textsuperscript{1}: Cook, \emph{Exercises in Old English} (1895), XII, XIII.37–43, printed pp. 22–28 (PDF 30–36); \emph{A First Book in Old English}, second edition (1895), especially §§55, 153–166, 183, 188, 194–197 and Reader XII, 174.7–175.17 (PDF 190–191). Individual line references count physical prose lines from the edition's printed marginal anchors. The supplied 144,000 expression is independently attested in Bosworth–Toller, main dictionary p. 1082, \textbf{þūsend} I(2)(a), citing \emph{Pastoral Care} 52; its acute length signs are represented by Cook-style macrons. The readings and original note numbers remain those of the selected editions.
% mdparagraph:P553:end

\end{document}
